%%% element catalog, generated on Tue May 19 10:30:56 CEST 2026

%%% WARNING %%%
%%% This file was automatically generated by SAPlugin,
%%% and will be overwritten next time.

\chapter{Element catalog}\label{sec:catalog}
% COMPONENTS
\section{Components}\label{sec:components}
\subsection{AlertManager}\label{comp:ComponentseHealthPlatformAlertManager}
	\begin{description}[noitemsep,nolistsep]
		\item[Responsibility:]~Manages the Quality of Response for all system notifications. It receives failure signals, determines the patient's risk level, and ensures the notification reaches the correct stakeholder (System Admin or Operator) within the mandated deadlines.
		\item[Super-components:]~\iconcomponent{}~\vpett{\nameref{comp:ComponentseHealthPlatform}}
		\item[Sub-components:]~None
		\item[Provided interfaces:]~\iconprovided{}~\vpett{\nameref{int:InterfacesAlertMgmt}}, \iconprovided{}~\vpett{\nameref{int:InterfacesBatteryAlertMgmt}}
		\item[Required interfaces:]~\iconrequired{}~\vpett{\nameref{int:InterfacesAlertPush}}, \iconrequired{}~\vpett{\nameref{int:InterfacesOtherDataMgmt}}
		\item[Deployed on:]~\faSquareO~\vpett{\nameref{node:DeploymentControlNode}}
		\item[Visible on diagrams:]~\cref{diag:Component:Contextdiagramfortheclientserverview,diag:Component:Primarydiagramoftheclientserverview,diag:Deployment:Primarydeploymentdiagram,diag:Interaction:Patientgatewayfails,diag:Interaction:av2CommunicationchannelbetweenPMSandPGgoesdownDetectionbytimeoutBackend}		
	\end{description}

\subsection{ClinicalJobCreator}\label{comp:ComponentseHealthPlatformClinicalJobCreator}
	\begin{description}[noitemsep,nolistsep]
		\item[Responsibility:]~The \vpett{\nameref{comp:ComponentseHealthPlatformClinicalJobCreator}} is responsible for constructing the \ref{data:DatatypesClinicalModelJob} objects and populating them with the necessary patient health data for execution by the \vpett{\nameref{comp:ComponentseHealthPlatformRiskEstimationProcessor}}s.
		\item[Super-components:]~\iconcomponent{}~\vpett{\nameref{comp:ComponentseHealthPlatform}}
		\item[Sub-components:]~None
		\item[Provided interfaces:]~\iconprovided{}~\vpett{\nameref{int:InterfacesClinicalJobMgmt}}
		\item[Required interfaces:]~\iconrequired{}~\vpett{\nameref{int:InterfacesFetchClinicalModels}}, \iconrequired{}~\vpett{\nameref{int:InterfacesOtherDataMgmt}}, \iconrequired{}~\vpett{\nameref{int:InterfacesSensorDataMgmt}}, \iconrequired{}~\vpett{\nameref{int:InterfacesThirdPartyDataMgmt}}
		\item[Deployed on:]~\faSquareO~\vpett{\nameref{node:DeploymentRiskEstimationManagementNodePilotDeployment}}, \faSquareO~\vpett{\nameref{node:DeploymentRiskEstimationManagmentNodeDevDeployment}}, \faSquareO~\vpett{\nameref{node:DeploymenteHealthPlatformRiskEstimationMgmtNode}}
		\item[Visible on diagrams:]~\cref{diag:Component:Primarydiagramoftheclientserverview,diag:Deployment:DevelopmentTestDeployment20patients3clinmodels,diag:Deployment:PilotDeployment200patients5clinmodels,diag:Deployment:Primarydeploymentdiagram,diag:Interaction:Riskestimationprocess,diag:Interaction:Thirdpartypolling}		
	\end{description}

\subsection{ClinicalModelCache}\label{comp:ComponentseHealthPlatformClinicalModelCache}
	\begin{description}[noitemsep,nolistsep]
		\item[Responsibility:]~The \vpett{\nameref{comp:ComponentseHealthPlatformClinicalModelCache}} is a read-through cache responsible for caching the \ref{data:DatatypesClinicalModel}s that should be evaluated for a certain patient and their configurations. This cache is located close to the \vpett{\nameref{comp:ComponentseHealthPlatformRiskEstimationProcessor}} and \vpett{\nameref{comp:ComponentseHealthPlatformRiskEstimationCombiner}} in order to improve the latency of the whole risk estimation flow. The items in the \vpett{\nameref{comp:ComponentseHealthPlatformClinicalModelCache}} do not expire over time, but should be invalided explicitly if needed.
Cached models remain available for the \vpett{\nameref{comp:ComponentseHealthPlatformRiskEstimationProcessor}} while updates are made in the \vpett{\nameref{comp:ComponentseHealthPlatformClinicalModelDB}}.
		\item[Super-components:]~\iconcomponent{}~\vpett{\nameref{comp:ComponentseHealthPlatform}}
		\item[Sub-components:]~None
		\item[Provided interfaces:]~\iconprovided{}~\vpett{\nameref{int:InterfacesClinicalModelCacheMgmt}}, \iconprovided{}~\vpett{\nameref{int:InterfacesFetchClinicalModels}}
		\item[Required interfaces:]~\iconrequired{}~\vpett{\nameref{int:InterfacesClinicalModelStorage}}, \iconrequired{}~\vpett{\nameref{int:InterfacesOtherDataMgmt}}
		\item[Deployed on:]~\faSquareO~\vpett{\nameref{node:DeploymentRiskEstimationManagementNodePilotDeployment}}, \faSquareO~\vpett{\nameref{node:DeploymentRiskEstimationManagmentNodeDevDeployment}}, \faSquareO~\vpett{\nameref{node:DeploymenteHealthPlatformRiskEstimationMgmtNode}}
		\item[Visible on diagrams:]~\cref{diag:Component:Contextdiagramfortheclientserverview,diag:Component:Primarydiagramoftheclientserverview,diag:Deployment:DevelopmentTestDeployment20patients3clinmodels,diag:Deployment:PilotDeployment200patients5clinmodels,diag:Deployment:Primarydeploymentdiagram,diag:Interaction:Riskestimationprocess}		
	\end{description}

\subsection{ClinicalModelDB}\label{comp:ComponentseHealthPlatformClinicalModelDB}
	\begin{description}[noitemsep,nolistsep]
		\item[Responsibility:]~The \vpett{\nameref{comp:ComponentseHealthPlatformClinicalModelDB}} stores all the \ref{data:DatatypesClinicalModel}s and the configurations of these \ref{data:DatatypesClinicalModel}s for the different patients separately from other data.
It only allows appending new \ref{data:DatatypesClinicalModel}s.
		\item[Super-components:]~\iconcomponent{}~\vpett{\nameref{comp:ComponentseHealthPlatform}}
		\item[Sub-components:]~None
		\item[Provided interfaces:]~\iconprovided{}~\vpett{\nameref{int:InterfacesClinicalModelStorage}}
		\item[Required interfaces:]~None
		\item[Deployed on:]~\faSquareO~\vpett{\nameref{node:DeploymentRiskEstimationManagmentNodeDevDeployment}}, \faSquareO~\vpett{\nameref{node:DeploymentModelStorage}}, \faSquareO~\vpett{\nameref{node:DeploymentmodelDBPilotDeployment}}
		\item[Visible on diagrams:]~\cref{diag:Component:Primarydiagramoftheclientserverview,diag:Deployment:DevelopmentTestDeployment20patients3clinmodels,diag:Deployment:PilotDeployment200patients5clinmodels,diag:Deployment:Primarydeploymentdiagram,diag:Interaction:Riskestimationprocess}		
	\end{description}

\subsection{Cronjob}\label{comp:ComponentseHealthPlatformCronjob}
	\begin{description}[noitemsep,nolistsep]
		\item[Responsibility:]~A component for running tasks at regular time intervals
		\item[Super-components:]~\iconcomponent{}~\vpett{\nameref{comp:ComponentseHealthPlatform}}
		\item[Sub-components:]~None
		\item[Provided interfaces:]~\iconprovided{}~\vpett{\nameref{int:InterfacesResetTimer}}
		\item[Required interfaces:]~\iconrequired{}~\vpett{\nameref{int:InterfacesAlertMgmt}}, \iconrequired{}~\vpett{\nameref{int:InterfacesCommuniationProblems}}
		\item[Deployed on:]~\faSquareO~\vpett{\nameref{node:DeploymentControlNode}}
		\item[Visible on diagrams:]~\cref{diag:Component:Contextdiagramfortheclientserverview,diag:Component:Primarydiagramoftheclientserverview,diag:Deployment:DevelopmentTestDeployment20patients3clinmodels,diag:Deployment:PilotDeployment200patients5clinmodels,diag:Deployment:Primarydeploymentdiagram,diag:Interaction:Incomingsensordata,diag:Interaction:av2CommunicationchannelbetweenPMSandPGgoesdownDetectionbytimeoutBackend}		
	\end{description}

\subsection{eHealth Platform}\label{comp:ComponentseHealthPlatform}
	\begin{description}[noitemsep,nolistsep]
		\item[Responsibility:]~The parent \vpett{\nameref{comp:ComponentseHealthPlatform}} component that represents the context boundary in the client-server view.
		\item[Super-components:]~None
		\item[Sub-components:]~\iconcomponent{}~\vpett{\nameref{comp:ComponentseHealthPlatformMLModelManager}}, \iconcomponent{}~\vpett{\nameref{comp:ComponentseHealthPlatformClinicalModelCache}}, \iconcomponent{}~\vpett{\nameref{comp:ComponentseHealthPlatformRiskEstimationCombiner}}, \iconcomponent{}~\vpett{\nameref{comp:ComponentseHealthPlatformRiskEstimationProcessor}}, \iconcomponent{}~\vpett{\nameref{comp:ComponentseHealthPlatformRiskEstimationScheduler}}, \iconcomponent{}~\vpett{\nameref{comp:ComponentseHealthPlatformClinicalJobCreator}}, \iconcomponent{}~\vpett{\nameref{comp:ComponentseHealthPlatformClinicalModelDB}}, \iconcomponent{}~\vpett{\nameref{comp:ComponentseHealthPlatformPatientDataStorage}}, \iconcomponent{}~\vpett{\nameref{comp:ComponentseHealthPlatformCronjob}}, \iconcomponent{}~\vpett{\nameref{comp:ComponentseHealthPlatformHIS}}, \iconcomponent{}~\vpett{\nameref{comp:ComponentseHealthPlatformHISManager}}, \iconcomponent{}~\vpett{\nameref{comp:ComponentseHealthPlatformAlertManager}}
		\item[Provided interfaces:]~\iconprovided{}~\vpett{\nameref{int:InterfacesBatteryAlertMgmt}}, \iconprovided{}~\vpett{\nameref{int:InterfacesClinicalModelCacheMgmt}}, \iconprovided{}~\vpett{\nameref{int:InterfacesLaunchRiskEstimation}}, \iconprovided{}~\vpett{\nameref{int:InterfacesResetTimer}}
		\item[Required interfaces:]~\iconrequired{}~\vpett{\nameref{int:InterfacesAlertPush}}, \iconrequired{}~\vpett{\nameref{int:InterfacesCommuniationProblems}}, \iconrequired{}~\vpett{\nameref{int:InterfacesKVStore}}, \iconrequired{}~\vpett{\nameref{int:InterfacesMLaaSService}}, \iconrequired{}~\vpett{\nameref{int:InterfacesThirdPartyApiInterface}}
		\item[Deployed on:]~\faSquareO~\vpett{\nameref{node:DeploymentRiskEstimationManagementNodePilotDeployment}}, \faSquareO~\vpett{\nameref{node:DeploymentRiskEstimationProcessorNodePilotDeployment}}, \faSquareO~\vpett{\nameref{node:DeploymentMLaaSConnPilotDeployment}}, \faSquareO~\vpett{\nameref{node:DeploymentRiskEstimationManagmentNodeDevDeployment}}, \faSquareO~\vpett{\nameref{node:DeploymentRiskEstimationProcessorNodeDevDeployment}}, \faSquareO~\vpett{\nameref{node:DeploymentDBNodePilotDeployment}}, \faSquareO~\vpett{\nameref{node:DeploymentPatientDB}}, \faSquareO~\vpett{\nameref{node:DeploymentModelStorage}}, \faSquareO~\vpett{\nameref{node:DeploymentmodelDBPilotDeployment}}, \faSquareO~\vpett{\nameref{node:DeploymentControlNode}}, \faSquareO~\vpett{\nameref{node:DeploymenteHealthPlatformMLaaSConnectorNode}}, \faSquareO~\vpett{\nameref{node:DeploymenteHealthPlatformRiskEstimationMgmtNode}}, \faSquareO~\vpett{\nameref{node:DeploymenteHealthPlatformRiskEstimationProcessorNode}}
		\item[Visible on diagrams:]~\cref{diag:Component:Contextdiagramfortheclientserverview,diag:Component:Contextdiagramforthedecompositionview}		
	\end{description}

\subsection{HIS}\label{comp:ComponentseHealthPlatformHIS}
	\begin{description}[noitemsep,nolistsep]
		\item[Responsibility:]~Hospital Information System
		\item[Super-components:]~\iconcomponent{}~\vpett{\nameref{comp:ComponentseHealthPlatform}}
		\item[Sub-components:]~None
		\item[Provided interfaces:]~\iconprovided{}~\vpett{\nameref{int:InterfaceshealthAPI}}
		\item[Required interfaces:]~None
		\item[Deployed on:]~\faSquareO~\vpett{\nameref{node:DeploymentControlNode}}
		\item[Visible on diagrams:]~\cref{diag:Component:Primarydiagramoftheclientserverview,diag:Deployment:DevelopmentTestDeployment20patients3clinmodels,diag:Deployment:PilotDeployment200patients5clinmodels,diag:Deployment:Primarydeploymentdiagram,diag:Interaction:Incomingsensordata}		
	\end{description}

\subsection{HIS Manager}\label{comp:ComponentseHealthPlatformHISManager}
	\begin{description}[noitemsep,nolistsep]
		\item[Responsibility:]~Manages calls to the \vpett{\nameref{comp:ComponentseHealthPlatformHIS}} and combines data recieved via healthAPI
		\item[Super-components:]~\iconcomponent{}~\vpett{\nameref{comp:ComponentseHealthPlatform}}
		\item[Sub-components:]~None
		\item[Provided interfaces:]~\iconprovided{}~\vpett{\nameref{int:InterfacesPatientRecordMgmt}}, \iconprovided{}~\vpett{\nameref{int:InterfacesSensorDataMgmt}}
		\item[Required interfaces:]~\iconrequired{}~\vpett{\nameref{int:InterfaceshealthAPI}}
		\item[Deployed on:]~\faSquareO~\vpett{\nameref{node:DeploymentControlNode}}
		\item[Visible on diagrams:]~\cref{diag:Component:Primarydiagramoftheclientserverview,diag:Deployment:DevelopmentTestDeployment20patients3clinmodels,diag:Deployment:PilotDeployment200patients5clinmodels,diag:Deployment:Primarydeploymentdiagram,diag:Interaction:Incomingsensordata}		
	\end{description}

\subsection{KVStorageService}\label{comp:ComponentsKVStorageService}
	\begin{description}[noitemsep,nolistsep]
		\item[Responsibility:]~Service for key-value storage of data.
		\item[Super-components:]~None
		\item[Sub-components:]~None
		\item[Provided interfaces:]~\iconprovided{}~\vpett{\nameref{int:InterfacesKVStore}}
		\item[Required interfaces:]~None
		\item[Deployed on:]~\faSquareO~\vpett{\nameref{node:DeploymentRiskEstimationProcessorNodePilotDeployment}}, \faSquareO~\vpett{\nameref{node:DeploymentRiskEstimationProcessorNodeDevDeployment}}, \faSquareO~\vpett{\nameref{node:DeploymenteHealthPlatformRiskEstimationProcessorNode}}
		\item[Visible on diagrams:]~\cref{diag:Component:Contextdiagramfortheclientserverview,diag:Deployment:DevelopmentTestDeployment20patients3clinmodels,diag:Deployment:PilotDeployment200patients5clinmodels,diag:Deployment:Primarydeploymentdiagram}		
	\end{description}

\subsection{MLaaS}\label{comp:ComponentsMLaaS}
	\begin{description}[noitemsep,nolistsep]
		\item[Responsibility:]~MachineLearning-as-a-Service cloud provider. This cloud service provides APIs for managing different \ref{data:DatatypesMLModel}s and exchanging \ref{data:DatatypesMLModel} updates with other parties.
		\item[Super-components:]~None
		\item[Sub-components:]~None
		\item[Provided interfaces:]~\iconprovided{}~\vpett{\nameref{int:InterfacesMLaaSService}}
		\item[Required interfaces:]~None
		\item[Deployed on:]~\faSquareO~\vpett{\nameref{node:DeploymentMLaaSCloud}}, \faSquareO~\vpett{\nameref{node:DeploymentMLaaSServicePilotDeployment}}, \faSquareO~\vpett{\nameref{node:DeploymentMLaaSServiceDevDeployment}}
		\item[Visible on diagrams:]~\cref{diag:Component:Contextdiagramfortheclientserverview,diag:Deployment:Contextdiagramforthedeploymentview,diag:Deployment:DevelopmentTestDeployment20patients3clinmodels,diag:Deployment:PilotDeployment200patients5clinmodels}		
	\end{description}

\subsection{MLModelManager}\label{comp:ComponentseHealthPlatformMLModelManager}
	\begin{description}[noitemsep,nolistsep]
		\item[Responsibility:]~This is component is responsible for managing the local \ref{data:DatatypesMLModel}s for making predictions for malignant events based on the incoming sensor data.
It keeps track of the available \ref{data:DatatypesMLModel}s that can be used and supports exchanging \ref{data:DatatypesMLModel} updates with other systems using the APIs of \vpett{\nameref{comp:ComponentsMLaaS}} providers.
		\item[Super-components:]~\iconcomponent{}~\vpett{\nameref{comp:ComponentseHealthPlatform}}
		\item[Sub-components:]~None
		\item[Provided interfaces:]~\iconprovided{}~\vpett{\nameref{int:InterfacesMLModelMgmt}}
		\item[Required interfaces:]~\iconrequired{}~\vpett{\nameref{int:InterfacesMLaaSService}}
		\item[Deployed on:]~\faSquareO~\vpett{\nameref{node:DeploymentMLaaSConnPilotDeployment}}, \faSquareO~\vpett{\nameref{node:DeploymentRiskEstimationManagmentNodeDevDeployment}}, \faSquareO~\vpett{\nameref{node:DeploymenteHealthPlatformMLaaSConnectorNode}}
		\item[Visible on diagrams:]~\cref{diag:Component:Contextdiagramfortheclientserverview,diag:Component:DecompositionoftheMLModelManager,diag:Component:Primarydiagramoftheclientserverview,diag:Deployment:Contextdiagramforthedeploymentview,diag:Deployment:DevelopmentTestDeployment20patients3clinmodels,diag:Deployment:PilotDeployment200patients5clinmodels,diag:Deployment:Primarydeploymentdiagram}		
	\end{description}

\subsection{Mobile OS}\label{comp:ComponentsMobileOS}
	\begin{description}[noitemsep,nolistsep]
		\item[Responsibility:]~Represents the mobile phone's operating system keeping track of battery, network connection, etc. 
		\item[Super-components:]~None
		\item[Sub-components:]~None
		\item[Provided interfaces:]~\iconprovided{}~\vpett{\nameref{int:InterfacesOSMgmt}}
		\item[Required interfaces:]~\iconrequired{}~\vpett{\nameref{int:InterfacesSystemEventListener}}
		\item[Deployed on:]~\faSquareO~\vpett{\nameref{node:DeploymentPatientDevices}}
		\item[Visible on diagrams:]~\cref{diag:Component:Contextdiagramfortheclientserverview,diag:Component:PatientGatewayAppComponentDiagram,diag:Deployment:Primarydeploymentdiagram,diag:Interaction:Patientgatewayfails}		
	\end{description}

\subsection{Operator Gateway Endpoint}\label{comp:ComponentsOperatorGatewayEndpoint}
	\begin{description}[noitemsep,nolistsep]
		\item[Responsibility:]~For handling incoming and outgoing requests from doctors, software testers and other operators
		\item[Super-components:]~None
		\item[Sub-components:]~None
		\item[Provided interfaces:]~\iconprovided{}~\vpett{\nameref{int:InterfacesAlertPush}}, \iconprovided{}~\vpett{\nameref{int:InterfacesCommuniationProblems}}
		\item[Required interfaces:]~\iconrequired{}~\vpett{\nameref{int:InterfacesClinicalModelCacheMgmt}}, \iconrequired{}~\vpett{\nameref{int:InterfacesGatewayRequests}}
		\item[Deployed on:]~\faSquareO~\vpett{\nameref{node:DeploymentControlNode}}
		\item[Visible on diagrams:]~\cref{diag:Component:Contextdiagramfortheclientserverview,diag:Component:Primarydiagramoftheclientserverview,diag:Deployment:DevelopmentTestDeployment20patients3clinmodels,diag:Deployment:PilotDeployment200patients5clinmodels,diag:Deployment:Primarydeploymentdiagram,diag:Interaction:Patientgatewayfails,diag:Interaction:Updategatewaypushwaittime,diag:Interaction:av2CommunicationchannelbetweenPMSandPGgoesdownDetectionbytimeoutBackend}		
	\end{description}

\subsection{Patient Gateway App}\label{comp:ComponentsPatientGatewayApp}
	\begin{description}[noitemsep,nolistsep]
		\item[Responsibility:]~The patient client side device
		\item[Super-components:]~None
		\item[Sub-components:]~None
		\item[Provided interfaces:]~\iconprovided{}~\vpett{\nameref{int:InterfacesPatientDataPush}}, \iconprovided{}~\vpett{\nameref{int:InterfacesSystemEventListener}}
		\item[Required interfaces:]~\iconrequired{}~\vpett{\nameref{int:InterfacesOSMgmt}}, \iconrequired{}~\vpett{\nameref{int:InterfacesPatientDataMgmt}}, \iconrequired{}~\vpett{\nameref{int:InterfacesSMSinterface}}
		\item[Deployed on:]~\faSquareO~\vpett{\nameref{node:DeploymentPatientDevices}}
		\item[Visible on diagrams:]~\cref{diag:Component:Contextdiagramfortheclientserverview,diag:Component:PatientGatewayAppComponentDiagram,diag:Component:Primarydiagramoftheclientserverview,diag:Deployment:Primarydeploymentdiagram,diag:Interaction:Patientgatewayfails,diag:Interaction:Storetokens,diag:Interaction:Updategatewaypushwaittime,diag:Interaction:av2CommunicationchannelbetweenPMSandPGgoesdownDetectionduringCommunicationAppSide}		
	\end{description}

\subsection{Patient Gateway Endpoint}\label{comp:ComponentsPatientGatewayEndpoint}
	\begin{description}[noitemsep,nolistsep]
		\item[Responsibility:]~The point where the pms connects to the gateway, this defines what the patient can do from his/her device and what to send to the gateway device
		\item[Super-components:]~None
		\item[Sub-components:]~None
		\item[Provided interfaces:]~\iconprovided{}~\vpett{\nameref{int:InterfacesGatewayRequests}}, \iconprovided{}~\vpett{\nameref{int:InterfacesPatientDataMgmt}}
		\item[Required interfaces:]~\iconrequired{}~\vpett{\nameref{int:InterfacesBatteryAlertMgmt}}, \iconrequired{}~\vpett{\nameref{int:InterfacesLaunchRiskEstimation}}, \iconrequired{}~\vpett{\nameref{int:InterfacesPatientDataPush}}, \iconrequired{}~\vpett{\nameref{int:InterfacesResetTimer}}, \iconrequired{}~\vpett{\nameref{int:InterfacesSensorDataMgmt}}, \iconrequired{}~\vpett{\nameref{int:InterfacesTokenMgmt}}
		\item[Deployed on:]~\faSquareO~\vpett{\nameref{node:DeploymentControlNode}}
		\item[Visible on diagrams:]~\cref{diag:Component:Contextdiagramfortheclientserverview,diag:Component:PatientGatewayAppComponentDiagram,diag:Component:Primarydiagramoftheclientserverview,diag:Deployment:DevelopmentTestDeployment20patients3clinmodels,diag:Deployment:PilotDeployment200patients5clinmodels,diag:Deployment:Primarydeploymentdiagram,diag:Interaction:Incomingsensordata,diag:Interaction:Patientgatewayfails,diag:Interaction:Storetokens,diag:Interaction:Updategatewaypushwaittime,diag:Interaction:av2CommunicationchannelbetweenPMSandPGgoesdownDetectionbytimeoutBackend,diag:Interaction:av2CommunicationchannelbetweenPMSandPGgoesdownDetectionduringCommunicationAppSide}		
	\end{description}

\subsection{PatientDataStorage}\label{comp:ComponentseHealthPlatformPatientDataStorage}
	\begin{description}[noitemsep,nolistsep]
		\item[Responsibility:]~The local patient db and interface to the db for storing third party patient data, and other stuff needed but not savable in the hospital informatioon system (\vpett{\nameref{comp:ComponentseHealthPlatformHIS}})
		\item[Super-components:]~\iconcomponent{}~\vpett{\nameref{comp:ComponentseHealthPlatform}}
		\item[Sub-components:]~None
		\item[Provided interfaces:]~\iconprovided{}~\vpett{\nameref{int:InterfacesOtherDataMgmt}}, \iconprovided{}~\vpett{\nameref{int:InterfacesThirdPartyDataMgmt}}
		\item[Required interfaces:]~\iconrequired{}~\vpett{\nameref{int:InterfacesThirdPartyApiInterface}}
		\item[Deployed on:]~\faSquareO~\vpett{\nameref{node:DeploymentDBNodePilotDeployment}}, \faSquareO~\vpett{\nameref{node:DeploymentPatientDB}}
		\item[Visible on diagrams:]~\cref{diag:Component:Contextdiagramfortheclientserverview,diag:Component:Primarydiagramoftheclientserverview,diag:Deployment:DevelopmentTestDeployment20patients3clinmodels,diag:Deployment:PilotDeployment200patients5clinmodels,diag:Deployment:Primarydeploymentdiagram,diag:Interaction:Incomingsensordata,diag:Interaction:Thirdpartypolling}		
	\end{description}

\subsection{RiskEstimationCombiner}\label{comp:ComponentseHealthPlatformRiskEstimationCombiner}
	\begin{description}[noitemsep,nolistsep]
		\item[Responsibility:]~The \vpett{\nameref{comp:ComponentseHealthPlatformRiskEstimationCombiner}} is responsible for combining the results of the \ref{data:DatatypesClinicalModelJob}s belonging to a patient risk estimation.
More precisely, the \vpett{\nameref{comp:ComponentseHealthPlatformRiskEstimationScheduler}} passes the set of the scheduled jobs for a risk estimation to the \vpett{\nameref{comp:ComponentseHealthPlatformRiskEstimationCombiner}} before scheduling them.
The \vpett{\nameref{comp:ComponentseHealthPlatformRiskEstimationCombiner}} then waits for all results to arrive, combines them, and propagates the new sensor data and the results of the risk estimation to the patient record if needed.
		\item[Super-components:]~\iconcomponent{}~\vpett{\nameref{comp:ComponentseHealthPlatform}}
		\item[Sub-components:]~None
		\item[Provided interfaces:]~\iconprovided{}~\vpett{\nameref{int:InterfacesJobMgmt}}, \iconprovided{}~\vpett{\nameref{int:InterfacesResults}}
		\item[Required interfaces:]~\iconrequired{}~\vpett{\nameref{int:InterfacesOtherDataMgmt}}, \iconrequired{}~\vpett{\nameref{int:InterfacesPatientRecordMgmt}}
		\item[Deployed on:]~\faSquareO~\vpett{\nameref{node:DeploymentRiskEstimationManagementNodePilotDeployment}}, \faSquareO~\vpett{\nameref{node:DeploymentRiskEstimationManagmentNodeDevDeployment}}, \faSquareO~\vpett{\nameref{node:DeploymenteHealthPlatformRiskEstimationMgmtNode}}
		\item[Visible on diagrams:]~\cref{diag:Component:Primarydiagramoftheclientserverview,diag:Deployment:DevelopmentTestDeployment20patients3clinmodels,diag:Deployment:PilotDeployment200patients5clinmodels,diag:Deployment:Primarydeploymentdiagram,diag:Interaction:CombiningClinicalModelJobresults,diag:Interaction:Incomingsensordata,diag:Interaction:Riskestimationprocess}		
	\end{description}

\subsection{RiskEstimationProcessor}\label{comp:ComponentseHealthPlatformRiskEstimationProcessor}
	\begin{description}[noitemsep,nolistsep]
		\item[Responsibility:]~The \vpett{\nameref{comp:ComponentseHealthPlatformRiskEstimationProcessor}} is responsible for computing \ref{data:DatatypesClinicalModelJob}s. The
\vpett{\nameref{comp:ComponentseHealthPlatformRiskEstimationProcessor}} fetches new \ref{data:DatatypesClinicalModelJob}s from the \vpett{\nameref{comp:ComponentseHealthPlatformRiskEstimationScheduler}}, uses the specified \ref{data:DatatypesMLModel} to calculate the risk prediction, and passes the result to the \vpett{\nameref{comp:ComponentseHealthPlatformRiskEstimationCombiner}}.
Multiple instances of the \vpett{\nameref{comp:ComponentseHealthPlatformRiskEstimationProcessor}} run in parallel to improve the throughput.
		\item[Super-components:]~\iconcomponent{}~\vpett{\nameref{comp:ComponentseHealthPlatform}}
		\item[Sub-components:]~None
		\item[Provided interfaces:]~None
		\item[Required interfaces:]~\iconrequired{}~\vpett{\nameref{int:InterfacesFetchJobs}}, \iconrequired{}~\vpett{\nameref{int:InterfacesKVStore}}, \iconrequired{}~\vpett{\nameref{int:InterfacesMLModelMgmt}}, \iconrequired{}~\vpett{\nameref{int:InterfacesResults}}
		\item[Deployed on:]~\faSquareO~\vpett{\nameref{node:DeploymentRiskEstimationProcessorNodePilotDeployment}}, \faSquareO~\vpett{\nameref{node:DeploymentRiskEstimationProcessorNodeDevDeployment}}, \faSquareO~\vpett{\nameref{node:DeploymenteHealthPlatformRiskEstimationProcessorNode}}
		\item[Visible on diagrams:]~\cref{diag:Component:Contextdiagramfortheclientserverview,diag:Component:DecompositionoftheRiskEstimationProcessor,diag:Component:Primarydiagramoftheclientserverview,diag:Deployment:DevelopmentTestDeployment20patients3clinmodels,diag:Deployment:PilotDeployment200patients5clinmodels,diag:Deployment:Primarydeploymentdiagram,diag:Interaction:CombiningClinicalModelJobresults}		
	\end{description}

\subsection{RiskEstimationScheduler}\label{comp:ComponentseHealthPlatformRiskEstimationScheduler}
	\begin{description}[noitemsep,nolistsep]
		\item[Responsibility:]~The \vpett{\nameref{comp:ComponentseHealthPlatformRiskEstimationScheduler}} is responsible for taking in new and scheduling requests for risk estimations.
It keeps track of the throughput of incoming jobs and changes the scheduling from first-in/first-out to dynamic priority earliest deadline first when going into overload modus.
		\item[Super-components:]~\iconcomponent{}~\vpett{\nameref{comp:ComponentseHealthPlatform}}
		\item[Sub-components:]~None
		\item[Provided interfaces:]~\iconprovided{}~\vpett{\nameref{int:InterfacesFetchJobs}}, \iconprovided{}~\vpett{\nameref{int:InterfacesLaunchRiskEstimation}}
		\item[Required interfaces:]~\iconrequired{}~\vpett{\nameref{int:InterfacesClinicalJobMgmt}}, \iconrequired{}~\vpett{\nameref{int:InterfacesFetchClinicalModels}}, \iconrequired{}~\vpett{\nameref{int:InterfacesJobMgmt}}, \iconrequired{}~\vpett{\nameref{int:InterfacesOtherDataMgmt}}, \iconrequired{}~\vpett{\nameref{int:InterfacesSensorDataMgmt}}
		\item[Deployed on:]~\faSquareO~\vpett{\nameref{node:DeploymentRiskEstimationManagementNodePilotDeployment}}, \faSquareO~\vpett{\nameref{node:DeploymentRiskEstimationManagmentNodeDevDeployment}}, \faSquareO~\vpett{\nameref{node:DeploymenteHealthPlatformRiskEstimationMgmtNode}}
		\item[Visible on diagrams:]~\cref{diag:Component:Contextdiagramfortheclientserverview,diag:Component:Primarydiagramoftheclientserverview,diag:Deployment:DevelopmentTestDeployment20patients3clinmodels,diag:Deployment:PilotDeployment200patients5clinmodels,diag:Deployment:Primarydeploymentdiagram,diag:Interaction:Incomingsensordata,diag:Interaction:Riskestimationprocess}		
	\end{description}

\subsection{SMS gateway}\label{comp:ComponentsSMSgateway}
	\begin{description}[noitemsep,nolistsep]
		\item[Responsibility:]~An external telecommunication system that provides an alternative, low-bandwidth communication channel via SMS.
		\item[Super-components:]~None
		\item[Sub-components:]~None
		\item[Provided interfaces:]~\iconprovided{}~\vpett{\nameref{int:InterfacesSMSinterface}}
		\item[Required interfaces:]~None
		\item[Deployed on:]~\faSquareO~\vpett{\nameref{node:DeploymenteHealthPlatformSMSProvider}}
		\item[Visible on diagrams:]~\cref{diag:Component:Contextdiagramfortheclientserverview,diag:Component:PatientGatewayAppComponentDiagram,diag:Deployment:Primarydeploymentdiagram,diag:Interaction:av2CommunicationchannelbetweenPMSandPGgoesdownDetectionduringCommunicationAppSide}		
	\end{description}

\subsection{ThirdPartyService}\label{comp:ComponentsThirdPartyService}
	\begin{description}[noitemsep,nolistsep]
		\item[Responsibility:]~the service that holds everything needed for the connection with thirdparty apps

		\item[Super-components:]~None
		\item[Sub-components:]~None
		\item[Provided interfaces:]~\iconprovided{}~\vpett{\nameref{int:InterfacesThirdPartyApiInterface}}, \iconprovided{}~\vpett{\nameref{int:InterfacesTokenMgmt}}
		\item[Required interfaces:]~None
		\item[Deployed on:]~\faSquareO~\vpett{\nameref{node:DeploymentDBNodePilotDeployment}}, \faSquareO~\vpett{\nameref{node:DeploymentPatientDB}}
		\item[Visible on diagrams:]~\cref{diag:Component:Contextdiagramfortheclientserverview,diag:Component:DecompositionThirdPartyApps,diag:Component:Primarydiagramoftheclientserverview,diag:Deployment:DevelopmentTestDeployment20patients3clinmodels,diag:Deployment:PilotDeployment200patients5clinmodels,diag:Deployment:Primarydeploymentdiagram,diag:Interaction:Storetokens,diag:Interaction:Thirdpartypolling}		
	\end{description}


% END COMPONENTS

% MODULES
\section{Modules}\label{sec:modules}
\subsection{App Endpoint}\label{mod:ComponentsPatientGatewayAppAppEndpoint}
	\begin{description}[noitemsep,nolistsep]
		\item[Responsibility:]~the component responsible for all communication from the app to the outside
		\item[Super-components:]~\iconcomponent{}~\vpett{\nameref{comp:ComponentsPatientGatewayApp}}
		\item[Super-modules:]~None
		\item[Sub-modules:]~None
		\item[Provided interfaces:]~\iconprovided{}~\vpett{\nameref{int:InterfacesPatientDataPush}}, \iconprovided{}~\vpett{\nameref{int:InterfacesSyncMgmt}}, \iconprovided{}~\vpett{\nameref{int:InterfacesSystemEventListener}}
		\item[Required interfaces:]~\iconrequired{}~\vpett{\nameref{int:InterfacesOSMgmt}}, \iconrequired{}~\vpett{\nameref{int:InterfacesPatientDataMgmt}}, \iconrequired{}~\vpett{\nameref{int:InterfacesShutDown}}
		\item[Visible on diagrams:]~\cref{diag:Component:PatientGatewayAppComponentDiagram,diag:Deployment:Primarydeploymentdiagram,diag:Interaction:DegradedMode,diag:Interaction:Exitdegradedmode,diag:Interaction:Ping}		
	\end{description}

\subsection{Buffer}\label{mod:ComponentsPatientGatewayAppBuffer}
	\begin{description}[noitemsep,nolistsep]
		\item[Responsibility:]~A buffer in the patient gateway app for storing cached/buffered sensor readings
		\item[Super-components:]~\iconcomponent{}~\vpett{\nameref{comp:ComponentsPatientGatewayApp}}
		\item[Super-modules:]~None
		\item[Sub-modules:]~None
		\item[Provided interfaces:]~\iconprovided{}~\vpett{\nameref{int:InterfacesBufferMgmt}}
		\item[Required interfaces:]~None
		\item[Visible on diagrams:]~\cref{diag:Component:PatientGatewayAppComponentDiagram,diag:Deployment:Primarydeploymentdiagram,diag:Interaction:DegradedMode,diag:Interaction:Exitdegradedmode}		
	\end{description}

\subsection{DataRiskAssessment}\label{mod:ComponentseHealthPlatformRiskEstimationProcessorDataRiskAssessment}
	\begin{description}[noitemsep,nolistsep]
		\item[Responsibility:]~This module performs the actual risk assessment by using the \ref{data:DatatypesMLModel} to make a prediction on the sensor data.
		\item[Super-components:]~\iconcomponent{}~\vpett{\nameref{comp:ComponentseHealthPlatformRiskEstimationProcessor}} $\triangleright$ \iconcomponent{}~\vpett{\nameref{comp:ComponentseHealthPlatform}}
		\item[Super-modules:]~None
		\item[Sub-modules:]~None
		\item[Provided interfaces:]~\iconprovided{}~\vpett{\nameref{int:InterfacesDataAnalysis}}
		\item[Required interfaces:]~\iconrequired{}~\vpett{\nameref{int:InterfacesKVStore}}, \iconrequired{}~\vpett{\nameref{int:InterfacesMLAPI}}
		\item[Visible on diagrams:]~\cref{diag:Component:Contextdiagramforthedecompositionview,diag:Component:DecompositionoftheRiskEstimationProcessor,diag:Interaction:Computeclinicalmodelresult}		
	\end{description}

\subsection{Google Fit API}\label{mod:ComponentsThirdPartyServiceThirdPartyAPIsGoogleFitAPI}
	\begin{description}[noitemsep,nolistsep]
		\item[Responsibility:]~The implementation for communicating with the Google fit service
		\item[Super-components:]~None
		\item[Super-modules:]~\iconcomponent{}~\vpett{\nameref{mod:ComponentsThirdPartyServiceThirdPartyAPIs}}
		\item[Sub-modules:]~None
		\item[Provided interfaces:]~\iconprovided{}~\vpett{\nameref{int:InterfacesThirdPartyAPIConnection}}
		\item[Required interfaces:]~None
		\item[Visible on diagrams:]~\cref{diag:Component:DecompositionThirdPartyApps,diag:Interaction:ThirdPartyNewData}		
	\end{description}

\subsection{JobMLModelSelector}\label{mod:ComponentseHealthPlatformRiskEstimationProcessorJobMLModelSelector}
	\begin{description}[noitemsep,nolistsep]
		\item[Responsibility:]~This module is responsible for determining the appropriate \ref{data:DatatypesMLModel} to use for making the prediction on the provided sensor data, as there could be multiple \ref{data:DatatypesMLModel} candidates for making predictions.
		\item[Super-components:]~\iconcomponent{}~\vpett{\nameref{comp:ComponentseHealthPlatformRiskEstimationProcessor}} $\triangleright$ \iconcomponent{}~\vpett{\nameref{comp:ComponentseHealthPlatform}}
		\item[Super-modules:]~None
		\item[Sub-modules:]~None
		\item[Provided interfaces:]~\iconprovided{}~\vpett{\nameref{int:InterfacesApplyClinicalModel}}
		\item[Required interfaces:]~None
		\item[Visible on diagrams:]~\cref{diag:Component:DecompositionoftheRiskEstimationProcessor,diag:Interaction:Computeclinicalmodelresult}		
	\end{description}

\subsection{JobProcessor}\label{mod:ComponentseHealthPlatformRiskEstimationProcessorJobProcessor}
	\begin{description}[noitemsep,nolistsep]
		\item[Responsibility:]~The \vpett{\nameref{mod:ComponentseHealthPlatformRiskEstimationProcessorJobProcessor}} module coordinates the execution of the \ref{data:DatatypesClinicalModelJob}s to ensure that the required \ref{data:DatatypesMLModel}s are fetched, and the prediction is performed with the appropriate \ref{data:DatatypesMLModel}.
		\item[Super-components:]~\iconcomponent{}~\vpett{\nameref{comp:ComponentseHealthPlatformRiskEstimationProcessor}} $\triangleright$ \iconcomponent{}~\vpett{\nameref{comp:ComponentseHealthPlatform}}
		\item[Super-modules:]~None
		\item[Sub-modules:]~None
		\item[Provided interfaces:]~None
		\item[Required interfaces:]~\iconrequired{}~\vpett{\nameref{int:InterfacesApplyClinicalModel}}, \iconrequired{}~\vpett{\nameref{int:InterfacesDataAnalysis}}, \iconrequired{}~\vpett{\nameref{int:InterfacesFetchJobs}}, \iconrequired{}~\vpett{\nameref{int:InterfacesFetchMLModel}}, \iconrequired{}~\vpett{\nameref{int:InterfacesResults}}
		\item[Visible on diagrams:]~\cref{diag:Component:DecompositionoftheRiskEstimationProcessor,diag:Interaction:Computeclinicalmodelresult}		
	\end{description}

\subsection{MLaaSLibrary}\label{mod:ModulesMLaaSLibrary}
	\begin{description}[noitemsep,nolistsep]
		\item[Responsibility:]~\vpett{\nameref{comp:ComponentsMLaaS}} development library for working with the \vpett{\nameref{comp:ComponentsMLaaS}} providers.
It supports storage \& retrieval of \ref{data:DatatypesMLModel}s, and exchanging \ref{data:DatatypesMLModelUpdate}s with \vpett{\nameref{comp:ComponentsMLaaS}} providers.
		\item[Super-components:]~None
		\item[Super-modules:]~None
		\item[Sub-modules:]~None
		\item[Provided interfaces:]~\iconprovided{}~\vpett{\nameref{int:InterfacesMLaaSAPI}}
		\item[Required interfaces:]~None
		\item[Visible on diagrams:]~\cref{diag:Component:Contextdiagramforthedecompositionview}		
	\end{description}

\subsection{MLLibrary}\label{mod:ModulesMLLibrary}
	\begin{description}[noitemsep,nolistsep]
		\item[Responsibility:]~Library for working with \ref{data:DatatypesMLModel}s locally and using them to make predictions based on the incoming sensor data.
The library can also process external updates to apply them to a specific \ref{data:DatatypesMLModel}.
		\item[Super-components:]~None
		\item[Super-modules:]~None
		\item[Sub-modules:]~None
		\item[Provided interfaces:]~\iconprovided{}~\vpett{\nameref{int:InterfacesMLAPI}}
		\item[Required interfaces:]~None
		\item[Visible on diagrams:]~\cref{diag:Component:Contextdiagramforthedecompositionview}		
	\end{description}

\subsection{MLModelRepository}\label{mod:ComponentseHealthPlatformMLModelManagerMLModelRepository}
	\begin{description}[noitemsep,nolistsep]
		\item[Responsibility:]~This module stores \ref{data:DatatypesMLModel}s locally and makes them available for use by the \vpett{\nameref{comp:ComponentseHealthPlatformRiskEstimationProcessor}}  when processing \ref{data:DatatypesClinicalModelJob}s.
		\item[Super-components:]~\iconcomponent{}~\vpett{\nameref{comp:ComponentseHealthPlatform}} $\triangleright$ \iconcomponent{}~\vpett{\nameref{comp:ComponentseHealthPlatformMLModelManager}}
		\item[Super-modules:]~None
		\item[Sub-modules:]~None
		\item[Provided interfaces:]~\iconprovided{}~\vpett{\nameref{int:InterfacesMLModelStorage}}
		\item[Required interfaces:]~None
		\item[Visible on diagrams:]~\cref{diag:Component:DecompositionoftheMLModelManager,diag:Interaction:MLModelSynchronization,diag:Interaction:ProcessingMLModelupdates}		
	\end{description}

\subsection{MLModelRetrieval}\label{mod:ComponentseHealthPlatformRiskEstimationProcessorMLModelRetrieval}
	\begin{description}[noitemsep,nolistsep]
		\item[Responsibility:]~This module is responsible for retrieving the requested \ref{data:DatatypesMLModel}s from the \vpett{\nameref{comp:ComponentseHealthPlatformMLModelManager}}.
		\item[Super-components:]~\iconcomponent{}~\vpett{\nameref{comp:ComponentseHealthPlatformRiskEstimationProcessor}} $\triangleright$ \iconcomponent{}~\vpett{\nameref{comp:ComponentseHealthPlatform}}
		\item[Super-modules:]~None
		\item[Sub-modules:]~None
		\item[Provided interfaces:]~\iconprovided{}~\vpett{\nameref{int:InterfacesFetchMLModel}}
		\item[Required interfaces:]~\iconrequired{}~\vpett{\nameref{int:InterfacesMLModelMgmt}}
		\item[Visible on diagrams:]~\cref{diag:Component:DecompositionoftheRiskEstimationProcessor,diag:Interaction:Computeclinicalmodelresult}		
	\end{description}

\subsection{MLModelStorageManager}\label{mod:ComponentseHealthPlatformMLModelManagerMLModelStorageManager}
	\begin{description}[noitemsep,nolistsep]
		\item[Responsibility:]~This module is responsible for managing the locally stored \ref{data:DatatypesMLModel}s, triggering the synchronization and retrieval of new \ref{data:DatatypesMLModel}s, and coordinating the updates to the existing \ref{data:DatatypesMLModel}s.
		\item[Super-components:]~\iconcomponent{}~\vpett{\nameref{comp:ComponentseHealthPlatform}} $\triangleright$ \iconcomponent{}~\vpett{\nameref{comp:ComponentseHealthPlatformMLModelManager}}
		\item[Super-modules:]~None
		\item[Sub-modules:]~None
		\item[Provided interfaces:]~\iconprovided{}~\vpett{\nameref{int:InterfacesMLModelMgmt}}
		\item[Required interfaces:]~\iconrequired{}~\vpett{\nameref{int:InterfacesMLModelStorage}}, \iconrequired{}~\vpett{\nameref{int:InterfacesMLModelSync}}
		\item[Visible on diagrams:]~\cref{diag:Component:DecompositionoftheMLModelManager,diag:Interaction:MLModelSynchronization,diag:Interaction:ProcessingMLModelupdates}		
	\end{description}

\subsection{MLModelUpdateProcessor}\label{mod:ComponentseHealthPlatformMLModelManagerMLModelUpdateProcessor}
	\begin{description}[noitemsep,nolistsep]
		\item[Responsibility:]~The MLModelUpdater processes and applies updates to the local \ref{data:DatatypesMLModel}s to prevent updates to these models from requiring to retrieve the full \ref{data:DatatypesMLModel}s.
		\item[Super-components:]~\iconcomponent{}~\vpett{\nameref{comp:ComponentseHealthPlatform}} $\triangleright$ \iconcomponent{}~\vpett{\nameref{comp:ComponentseHealthPlatformMLModelManager}}
		\item[Super-modules:]~None
		\item[Sub-modules:]~None
		\item[Provided interfaces:]~\iconprovided{}~\vpett{\nameref{int:InterfacesMLModelSync}}
		\item[Required interfaces:]~\iconrequired{}~\vpett{\nameref{int:InterfacesMLaaSAPI}}, \iconrequired{}~\vpett{\nameref{int:InterfacesMLaaSService}}, \iconrequired{}~\vpett{\nameref{int:InterfacesMLModelStorage}}, \iconrequired{}~\vpett{\nameref{int:InterfacesMLModelUpdateMgmt}}
		\item[Visible on diagrams:]~\cref{diag:Component:Contextdiagramforthedecompositionview,diag:Component:DecompositionoftheMLModelManager,diag:Interaction:MLModelSynchronization,diag:Interaction:ProcessingMLModelupdates}		
	\end{description}

\subsection{ModelUpdater}\label{mod:ComponentseHealthPlatformMLModelManagerModelUpdater}
	\begin{description}[noitemsep,nolistsep]
		\item[Responsibility:]~This module is responsible for processing \ref{data:DatatypesMLModelUpdate}s and applying them to the \ref{data:DatatypesMLModel}.
		\item[Super-components:]~\iconcomponent{}~\vpett{\nameref{comp:ComponentseHealthPlatform}} $\triangleright$ \iconcomponent{}~\vpett{\nameref{comp:ComponentseHealthPlatformMLModelManager}}
		\item[Super-modules:]~None
		\item[Sub-modules:]~None
		\item[Provided interfaces:]~\iconprovided{}~\vpett{\nameref{int:InterfacesMLModelUpdateMgmt}}
		\item[Required interfaces:]~\iconrequired{}~\vpett{\nameref{int:InterfacesMLAPI}}
		\item[Visible on diagrams:]~\cref{diag:Component:Contextdiagramforthedecompositionview,diag:Component:DecompositionoftheMLModelManager,diag:Interaction:ProcessingMLModelupdates}		
	\end{description}

\subsection{Strava API}\label{mod:ComponentsThirdPartyServiceThirdPartyAPIsStravaAPI}
	\begin{description}[noitemsep,nolistsep]
		\item[Responsibility:]~Api for communicating with the strava service
		\item[Super-components:]~None
		\item[Super-modules:]~\iconcomponent{}~\vpett{\nameref{mod:ComponentsThirdPartyServiceThirdPartyAPIs}}
		\item[Sub-modules:]~None
		\item[Provided interfaces:]~\iconprovided{}~\vpett{\nameref{int:InterfacesThirdPartyAPIConnection}}
		\item[Required interfaces:]~None
		\item[Visible on diagrams:]~\cref{diag:Component:DecompositionThirdPartyApps,diag:Interaction:ThirdPartyNewData}		
	\end{description}

\subsection{SyncManager}\label{mod:ComponentsPatientGatewayAppSyncManager}
	\begin{description}[noitemsep,nolistsep]
		\item[Responsibility:]~This component manages the state of the communication channel on the gateway. It is responsible for switching between "Normal" and "Degraded" modes.
		\item[Super-components:]~\iconcomponent{}~\vpett{\nameref{comp:ComponentsPatientGatewayApp}}
		\item[Super-modules:]~None
		\item[Sub-modules:]~None
		\item[Provided interfaces:]~\iconprovided{}~\vpett{\nameref{int:InterfacesShutDown}}
		\item[Required interfaces:]~\iconrequired{}~\vpett{\nameref{int:InterfacesBufferMgmt}}, \iconrequired{}~\vpett{\nameref{int:InterfacesSMSinterface}}, \iconrequired{}~\vpett{\nameref{int:InterfacesSyncMgmt}}
		\item[Visible on diagrams:]~\cref{diag:Component:PatientGatewayAppComponentDiagram,diag:Deployment:Primarydeploymentdiagram,diag:Interaction:DegradedMode,diag:Interaction:Exitdegradedmode,diag:Interaction:Ping}		
	\end{description}

\subsection{ThirdParty API's}\label{mod:ComponentsThirdPartyServiceThirdPartyAPIs}
	\begin{description}[noitemsep,nolistsep]
		\item[Responsibility:]~a component that can be implemented for every third party service that needs to be communicated with, if a new service needs to be available it needs to be implemented as a subclass of this class
		\item[Super-components:]~\iconcomponent{}~\vpett{\nameref{comp:ComponentsThirdPartyService}}
		\item[Super-modules:]~None
		\item[Sub-modules:]~\iconcomponent{}~\vpett{\nameref{mod:ComponentsThirdPartyServiceThirdPartyAPIsStravaAPI}}, \iconcomponent{}~\vpett{\nameref{mod:ComponentsThirdPartyServiceThirdPartyAPIsGoogleFitAPI}}
		\item[Provided interfaces:]~\iconprovided{}~\vpett{\nameref{int:InterfacesThirdPartyAPIConnection}}
		\item[Required interfaces:]~None
		\item[Visible on diagrams:]~\cref{diag:Component:DecompositionThirdPartyApps}		
	\end{description}

\subsection{ThirdPartySyncService}\label{mod:ComponentsThirdPartyServiceThirdPartySyncService}
	\begin{description}[noitemsep,nolistsep]
		\item[Responsibility:]~component for communicating with every third party service
		\item[Super-components:]~\iconcomponent{}~\vpett{\nameref{comp:ComponentsThirdPartyService}}
		\item[Super-modules:]~None
		\item[Sub-modules:]~None
		\item[Provided interfaces:]~\iconprovided{}~\vpett{\nameref{int:InterfacesThirdPartyApiInterface}}
		\item[Required interfaces:]~\iconrequired{}~\vpett{\nameref{int:InterfacesSearchToken}}, \iconrequired{}~\vpett{\nameref{int:InterfacesThirdPartyAPIConnection}}, \iconrequired{}~\vpett{\nameref{int:InterfacesTokenMgmt}}
		\item[Visible on diagrams:]~\cref{diag:Component:DecompositionThirdPartyApps,diag:Interaction:ThirdPartyNewData}		
	\end{description}

\subsection{TokenStorage}\label{mod:ComponentsThirdPartyServiceTokenStorage}
	\begin{description}[noitemsep,nolistsep]
		\item[Responsibility:]~a database to store tokens per patientId for third party apps.
		\item[Super-components:]~\iconcomponent{}~\vpett{\nameref{comp:ComponentsThirdPartyService}}
		\item[Super-modules:]~None
		\item[Sub-modules:]~None
		\item[Provided interfaces:]~\iconprovided{}~\vpett{\nameref{int:InterfacesSearchToken}}, \iconprovided{}~\vpett{\nameref{int:InterfacesTokenMgmt}}
		\item[Required interfaces:]~None
		\item[Visible on diagrams:]~\cref{diag:Component:DecompositionThirdPartyApps,diag:Interaction:ThirdPartyNewData}		
	\end{description}

\pdfinfo{/SAPlugin (v10.0.0)}
% END MODULES

% INTERFACES
\section{Interfaces} \label{sec:interfaces}
  %%%%%%%% AlertMgmt
  \subsection{AlertMgmt}\label{int:InterfacesAlertMgmt}
    \begin{description}[noitemsep,nolistsep]
      \item[Provided by:] \iconcomponent{}~\vpett{\nameref{comp:ComponentseHealthPlatformAlertManager}}
      \item[Required by:] \iconcomponent{}~\vpett{\nameref{comp:ComponentseHealthPlatformCronjob}} 
      \item[Operations:] ~
    \begin{itemize}[noitemsep,nolistsep,leftmargin=-.25cm]
      \item \textsf{void raiseAlert()}
        \begin{itemize}[noitemsep,nolistsep]
           \item Effect: Contact Operators and stakeholders in case of problem in system such as communication channel failure. 
           \item Sequence Diagrams: \cref{diag:Interaction:av2CommunicationchannelbetweenPMSandPGgoesdownDetectionbytimeoutBackend}
        \end{itemize}
    \end{itemize}
      \item[Diagrams:] None
    \end{description}

  %%%%%%%% AlertPush
  \subsection{AlertPush}\label{int:InterfacesAlertPush}
    \begin{description}[noitemsep,nolistsep]
      \item[Provided by:] \iconcomponent{}~\vpett{\nameref{comp:ComponentsOperatorGatewayEndpoint}}
      \item[Required by:] \iconcomponent{}~\vpett{\nameref{comp:ComponentseHealthPlatformAlertManager}}, \iconcomponent{}~\vpett{\nameref{comp:ComponentseHealthPlatform}} 
      \item[Operations:] ~
    \begin{itemize}[noitemsep,nolistsep,leftmargin=-.25cm]
      \item \textsf{void alertEmergencyContact(\ref{data:DatatypesPatientId} patientId)}
        \begin{itemize}[noitemsep,nolistsep]
           \item Effect: alert emergency contacts of the patient 
           \item Sequence Diagrams: \cref{diag:Interaction:Patientgatewayfails}
        \end{itemize}
      \item \textsf{void contactTelecom(\ref{data:DatatypesPatientId} patient, {\colorbox{red!30}{\underline{void}}} riskLevel)}
        \begin{itemize}[noitemsep,nolistsep]
           \item Effect: The system administrator needs to contact the telecommunication operator when the communication channel is down. 
           \item Sequence Diagrams: \cref{diag:Interaction:av2CommunicationchannelbetweenPMSandPGgoesdownDetectionbytimeoutBackend}
        \end{itemize}
    \end{itemize}
      \item[Diagrams:] None
    \end{description}

  %%%%%%%% ApplyClinicalModel
  \subsection{ApplyClinicalModel}\label{int:InterfacesApplyClinicalModel}
    \begin{description}[noitemsep,nolistsep]
      \item[Provided by:] \iconcomponent{}~\vpett{\nameref{mod:ComponentseHealthPlatformRiskEstimationProcessorJobMLModelSelector}}
      \item[Required by:] \iconcomponent{}~\vpett{\nameref{mod:ComponentseHealthPlatformRiskEstimationProcessorJobProcessor}} 
      \item[Operations:] ~
    \begin{itemize}[noitemsep,nolistsep,leftmargin=-.25cm]
      \item \textsf{String determineMLModelForClinicalModel(\ref{data:DatatypesClinicalModel} clinicalModel, \ref{data:DatatypesClinicalModelConfiguration} clinicalModelConfiguration)}
        \begin{itemize}[noitemsep,nolistsep]
           \item Effect: Apply the \ref{data:DatatypesClinicalModel} to obtain the most appropriate \ref{data:DatatypesMLModel} to use for the calculation. 
           \item Returns: \ref{data:Datatypesorghl7fhirr5modelIdentifier} of the \ref{data:DatatypesMLModel} to use.
           \item Sequence Diagrams: \cref{diag:Interaction:Computeclinicalmodelresult}
        \end{itemize}
    \end{itemize}
      \item[Diagrams:] None
    \end{description}

  %%%%%%%% BatteryAlertMgmt
  \subsection{BatteryAlertMgmt}\label{int:InterfacesBatteryAlertMgmt}
    \begin{description}[noitemsep,nolistsep]
      \item[Provided by:] \iconcomponent{}~\vpett{\nameref{comp:ComponentseHealthPlatformAlertManager}}, \iconcomponent{}~\vpett{\nameref{comp:ComponentseHealthPlatform}}
      \item[Required by:] \iconcomponent{}~\vpett{\nameref{comp:ComponentsPatientGatewayEndpoint}} 
      \item[Operations:] ~
    \begin{itemize}[noitemsep,nolistsep,leftmargin=-.25cm]
      \item \textsf{void sendBatteryAlert(\ref{data:DatatypesPatientId} patient)}
        \begin{itemize}[noitemsep,nolistsep]
           \item Effect: give the battery allert to the alert manager
 
           \item Sequence Diagrams: \cref{diag:Interaction:Patientgatewayfails}
        \end{itemize}
    \end{itemize}
      \item[Diagrams:] None
    \end{description}

  %%%%%%%% BufferMgmt
  \subsection{BufferMgmt}\label{int:InterfacesBufferMgmt}
    \begin{description}[noitemsep,nolistsep]
      \item[Provided by:] \iconcomponent{}~\vpett{\nameref{mod:ComponentsPatientGatewayAppBuffer}}
      \item[Required by:] \iconcomponent{}~\vpett{\nameref{mod:ComponentsPatientGatewayAppSyncManager}} 
           \item[Description]: to store/retrieve data in buffer in case of degraded mode.
      \item[Operations:] ~
    \begin{itemize}[noitemsep,nolistsep,leftmargin=-.25cm]
      \item \textsf{\ref{data:DatatypesSensorDataPackage} getNextChunk()}
        \begin{itemize}[noitemsep,nolistsep]
           \item Effect: get the next sensordatapackage in the buffer 
           \item Returns: the next sensor datapackage in the buffer
           \item Sequence Diagrams: \cref{diag:Interaction:Exitdegradedmode}
        \end{itemize}
      \item \textsf{boolean storeData(\ref{data:DatatypesSensorDataPackage} parameter)}
        \begin{itemize}[noitemsep,nolistsep]
           \item Effect: Called by \vpett{\nameref{mod:ComponentsPatientGatewayAppSyncManager}} to store data in case of degraded mode as backup, later retrieved to send to PMS backend. Returns false if the buffer was full and had to be overwritten. 
           \item Returns: Function returns true if buffer was not full when writing sensordata to it and false otherwise.
           \item Sequence Diagrams: \cref{diag:Interaction:DegradedMode}
        \end{itemize}
    \end{itemize}
      \item[Diagrams:] None
    \end{description}

  %%%%%%%% ClinicalJobMgmt
  \subsection{ClinicalJobMgmt}\label{int:InterfacesClinicalJobMgmt}
    \begin{description}[noitemsep,nolistsep]
      \item[Provided by:] \iconcomponent{}~\vpett{\nameref{comp:ComponentseHealthPlatformClinicalJobCreator}}
      \item[Required by:] \iconcomponent{}~\vpett{\nameref{comp:ComponentseHealthPlatformRiskEstimationScheduler}} 
      \item[Operations:] ~
    \begin{itemize}[noitemsep,nolistsep,leftmargin=-.25cm]
      \item \textsf{List\textless{}Tuple\textless{}\ref{data:DatatypesClinicalModelJob}, Double\textgreater{}\textgreater{} createClinicalModelJobsForPatient(\ref{data:DatatypesPatientId} patientId, Tuple\textless{}\ref{data:DatatypesSensorDataPackage}, \ref{data:DatatypesTimestamp}\textgreater{} sensorData)}
        \begin{itemize}[noitemsep,nolistsep]
           \item Effect: Construct the list of \ref{data:DatatypesClinicalModelJob}s for a specified patient. These jobs are populated with the necessary data for their execution by the \vpett{\nameref{comp:ComponentseHealthPlatformRiskEstimationProcessor}}. 
           \item Returns: A list of tuples containing the \ref{data:DatatypesClinicalModelJob} and its weight.
           \item Sequence Diagrams: \cref{diag:Interaction:Riskestimationprocess}
        \end{itemize}
    \end{itemize}
      \item[Diagrams:] None
    \end{description}

  %%%%%%%% ClinicalModelCacheMgmt
  \subsection{ClinicalModelCacheMgmt}\label{int:InterfacesClinicalModelCacheMgmt}
    \begin{description}[noitemsep,nolistsep]
      \item[Provided by:] \iconcomponent{}~\vpett{\nameref{comp:ComponentseHealthPlatformClinicalModelCache}}, \iconcomponent{}~\vpett{\nameref{comp:ComponentseHealthPlatform}}
      \item[Required by:] \iconcomponent{}~\vpett{\nameref{comp:ComponentsOperatorGatewayEndpoint}} 
      \item[Operations:] ~
    \begin{itemize}[noitemsep,nolistsep,leftmargin=-.25cm]
      \item \textsf{void invalidateCacheEntries(\ref{data:DatatypesPatientId} patientId)}
        \begin{itemize}[noitemsep,nolistsep]
           \item Effect: The \vpett{\nameref{comp:ComponentseHealthPlatformClinicalModelCache}} will invalidate (i.e., remove) all items in its cache for the patient identified by patientId. If the cache does not contain any items for this patient,
nothing is changed. After invalidating the cached items for a certain patient, the next request for them will lead to fetching them from the database and storing them in the cache
again. 
           \item Sequence Diagrams: None
        \end{itemize}
    \end{itemize}
      \item[Diagrams:] None
    \end{description}

  %%%%%%%% ClinicalModelStorage
  \subsection{ClinicalModelStorage}\label{int:InterfacesClinicalModelStorage}
    \begin{description}[noitemsep,nolistsep]
      \item[Provided by:] \iconcomponent{}~\vpett{\nameref{comp:ComponentseHealthPlatformClinicalModelDB}}
      \item[Required by:] \iconcomponent{}~\vpett{\nameref{comp:ComponentseHealthPlatformClinicalModelCache}} 
      \item[Operations:] ~
    \begin{itemize}[noitemsep,nolistsep,leftmargin=-.25cm]
      \item \textsf{Map\textless{}\ref{data:DatatypesClinicalModel}, Tuple\textless{}\ref{data:DatatypesClinicalModelConfiguration}, Double\textgreater{}\textgreater{} getAssignedClinicalModelsAndConfigForPatient(\ref{data:DatatypesPatientId} patientId) throws \ref{ex:ExceptionsNoSuchPatientException}}
        \begin{itemize}[noitemsep,nolistsep]
           \item Effect: Fetch and return the ids of the \ref{data:DatatypesClinicalModel}s assigned to the patient identified by patientId and their configurations for this patient. The configurations are returned as a \ref{data:DatatypesClinicalModelConfiguration} containing a map of configuration parameters and their value. 
           \item Returns: A map with for every applicable \ref{data:DatatypesClinicalModel}, a tuple containing the corresponding \ref{data:DatatypesClinicalModelConfiguration} for the specified patient and the corresponding weight factors used for combining the results.
           \item Sequence Diagrams: \cref{diag:Interaction:Riskestimationprocess}
        \end{itemize}
    \end{itemize}
      \item[Diagrams:] None
    \end{description}

  %%%%%%%% CommuniationProblems
  \subsection{CommuniationProblems}\label{int:InterfacesCommuniationProblems}
    \begin{description}[noitemsep,nolistsep]
      \item[Provided by:] \iconcomponent{}~\vpett{\nameref{comp:ComponentsOperatorGatewayEndpoint}}
      \item[Required by:] \iconcomponent{}~\vpett{\nameref{comp:ComponentseHealthPlatformCronjob}}, \iconcomponent{}~\vpett{\nameref{comp:ComponentseHealthPlatform}} 
      \item[Operations:] ~
    \begin{itemize}[noitemsep,nolistsep,leftmargin=-.25cm]
      \item \textsf{void recieveCommunicationProblemAlert(\ref{data:DatatypesPatientId} patientId)}
        \begin{itemize}[noitemsep,nolistsep]
           \item Effect: When the patient gateway endpoint does not recieve the sensordata a alert is send to the operator endpoint. 
           \item Sequence Diagrams: None
        \end{itemize}
    \end{itemize}
      \item[Diagrams:] None
    \end{description}

  %%%%%%%% DataAnalysis
  \subsection{DataAnalysis}\label{int:InterfacesDataAnalysis}
    \begin{description}[noitemsep,nolistsep]
      \item[Provided by:] \iconcomponent{}~\vpett{\nameref{mod:ComponentseHealthPlatformRiskEstimationProcessorDataRiskAssessment}}
      \item[Required by:] \iconcomponent{}~\vpett{\nameref{mod:ComponentseHealthPlatformRiskEstimationProcessorJobProcessor}} 
      \item[Operations:] ~
    \begin{itemize}[noitemsep,nolistsep,leftmargin=-.25cm]
      \item \textsf{void performModelPrediction(\ref{data:DatatypesClinicalModelJob} clinicalModelJob, \ref{data:DatatypesMLModel} model)}
        \begin{itemize}[noitemsep,nolistsep]
           \item Effect: Perform a prediction using the provided \ref{data:DatatypesMLModel} and the data provided in the clinicalModelJob definition.
Afterwards the results are submitted for combination by the \vpett{\nameref{comp:ComponentseHealthPlatformRiskEstimationCombiner}}. 
           \item Sequence Diagrams: \cref{diag:Interaction:Computeclinicalmodelresult}
        \end{itemize}
    \end{itemize}
      \item[Diagrams:] None
    \end{description}

  %%%%%%%% FetchClinicalModels
  \subsection{FetchClinicalModels}\label{int:InterfacesFetchClinicalModels}
    \begin{description}[noitemsep,nolistsep]
      \item[Provided by:] \iconcomponent{}~\vpett{\nameref{comp:ComponentseHealthPlatformClinicalModelCache}}
      \item[Required by:] \iconcomponent{}~\vpett{\nameref{comp:ComponentseHealthPlatformClinicalJobCreator}}, \iconcomponent{}~\vpett{\nameref{comp:ComponentseHealthPlatformRiskEstimationScheduler}} 
      \item[Operations:] ~
    \begin{itemize}[noitemsep,nolistsep,leftmargin=-.25cm]
      \item \textsf{Map\textless{}\ref{data:DatatypesClinicalModel}, Tuple\textless{}\ref{data:DatatypesClinicalModelConfiguration}, Double\textgreater{}\textgreater{} getAssignedClinicalModelsAndConfigForPatient(\ref{data:DatatypesPatientId} patientId)}
        \begin{itemize}[noitemsep,nolistsep]
           \item Effect: The \vpett{\nameref{comp:ComponentseHealthPlatformClinicalModelCache}} will return the clinical models assigned to the patient identified by patientId and their configurations for this patient. The configurations are
a map of configuration parameters and their value. If the cache contains data for the patient with given id, the data from the cache is returned. If not, the \vpett{\nameref{comp:ComponentseHealthPlatformClinicalModelCache}} will
fetch all clinical models assigned to the patient with given patientId and their configurations for this patient, store these in the cache and return them. 
           \item Returns: A map with for every applicable \ref{data:DatatypesClinicalModel}, a tuple containing the corresponding \ref{data:DatatypesClinicalModelConfiguration} for the specified patient and the corresponding weight factors used for combining the results.
           \item Sequence Diagrams: \cref{diag:Interaction:Riskestimationprocess}
        \end{itemize}
    \end{itemize}
      \item[Diagrams:] None
    \end{description}

  %%%%%%%% FetchJobs
  \subsection{FetchJobs}\label{int:InterfacesFetchJobs}
    \begin{description}[noitemsep,nolistsep]
      \item[Provided by:] \iconcomponent{}~\vpett{\nameref{comp:ComponentseHealthPlatformRiskEstimationScheduler}}
      \item[Required by:] \iconcomponent{}~\vpett{\nameref{mod:ComponentseHealthPlatformRiskEstimationProcessorJobProcessor}}, \iconcomponent{}~\vpett{\nameref{comp:ComponentseHealthPlatformRiskEstimationProcessor}} 
      \item[Operations:] ~
    \begin{itemize}[noitemsep,nolistsep,leftmargin=-.25cm]
      \item \textsf{\ref{data:DatatypesClinicalModelJob} getNextJob()}
        \begin{itemize}[noitemsep,nolistsep]
           \item Effect: Returns the next clinical model computation job that must be performed (i.e., the first job in the queue). 
           \item Returns: The next clinical model computation job that must be performed.
           \item Sequence Diagrams: \cref{diag:Interaction:Computeclinicalmodelresult}
        \end{itemize}
    \end{itemize}
      \item[Diagrams:] None
    \end{description}

  %%%%%%%% FetchMLModel
  \subsection{FetchMLModel}\label{int:InterfacesFetchMLModel}
    \begin{description}[noitemsep,nolistsep]
      \item[Provided by:] \iconcomponent{}~\vpett{\nameref{mod:ComponentseHealthPlatformRiskEstimationProcessorMLModelRetrieval}}
      \item[Required by:] \iconcomponent{}~\vpett{\nameref{mod:ComponentseHealthPlatformRiskEstimationProcessorJobProcessor}} 
      \item[Operations:] ~
    \begin{itemize}[noitemsep,nolistsep,leftmargin=-.25cm]
      \item \textsf{\ref{data:DatatypesMLModel} fetchModel(String modelId) throws \ref{ex:ExceptionsNoSuchMLModelException}}
        \begin{itemize}[noitemsep,nolistsep]
           \item Effect: Fetch the \ref{data:DatatypesMLModel} corresponding with the provided modelId. 
           \item Returns: The requested \ref{data:DatatypesMLModel}.
           \item Sequence Diagrams: \cref{diag:Interaction:Computeclinicalmodelresult}
        \end{itemize}
    \end{itemize}
      \item[Diagrams:] None
    \end{description}

  %%%%%%%% GatewayRequests
  \subsection{GatewayRequests}\label{int:InterfacesGatewayRequests}
    \begin{description}[noitemsep,nolistsep]
      \item[Provided by:] \iconcomponent{}~\vpett{\nameref{comp:ComponentsPatientGatewayEndpoint}}
      \item[Required by:] \iconcomponent{}~\vpett{\nameref{comp:ComponentsOperatorGatewayEndpoint}} 
      \item[Operations:] ~
    \begin{itemize}[noitemsep,nolistsep,leftmargin=-.25cm]
      \item \textsf{boolean changeTimings(Map\textless{}String, int\textgreater{} newTimings)}
        \begin{itemize}[noitemsep,nolistsep]
           \item Effect: change the timings specific sensor data should be sent over
 
           \item Returns: succes or not
           \item Sequence Diagrams: \cref{diag:Interaction:Updategatewaypushwaittime}
        \end{itemize}
      \item \textsf{\ref{data:DatatypesSensorDataPackage} getNewSensorData()}
        \begin{itemize}[noitemsep,nolistsep]
           \item Effect: get the new sensor data on the spot  
           \item Returns: the sesnordata of that moment
           \item Sequence Diagrams: None
        \end{itemize}
    \end{itemize}
      \item[Diagrams:] None
    \end{description}

  %%%%%%%% healthAPI
  \subsection{healthAPI}\label{int:InterfaceshealthAPI}
    \begin{description}[noitemsep,nolistsep]
      \item[Provided by:] \iconcomponent{}~\vpett{\nameref{comp:ComponentseHealthPlatformHIS}}
      \item[Required by:] \iconcomponent{}~\vpett{\nameref{comp:ComponentseHealthPlatformHISManager}} 
           \item[Description]: This interface provides a reduced set of simplified operations that are sufficient in the context of the patient monitoring service that is being developed.
It is derived from the HL7 FHIR standard for healthcare interoperability, based on the HAPI FHIR implementation of the FHIR specification in Java. The relevant OpenAPIs are described at: \url{https://hapi.fhir.org/baseR4/swagger-ui/}, while the FHIR standard is available at: \url{https://hl7.org/fhir/}.
The descriptions from the data types used in this interface follow directly from the API documentation (\url{https://hapifhir.io/hapi-fhir/apidocs/hapi-fhir-structures-r5/org/hl7/fhir/r5/model/package-summary.html}).
      \item[Operations:] ~
    \begin{itemize}[noitemsep,nolistsep,leftmargin=-.25cm]
      \item \textsf{void deleteObservation(string id)}
        \begin{itemize}[noitemsep,nolistsep]
           \item Effect: Delete an \ref{data:Datatypesorghl7fhirr5modelObservation} instance with the specified id. 
           \item Sequence Diagrams: None
        \end{itemize}
      \item \textsf{void deletePatient(string id)}
        \begin{itemize}[noitemsep,nolistsep]
           \item Effect: Delete a \ref{data:Datatypesorghl7fhirr5modelPatient} instance with the specfied id. 
           \item Sequence Diagrams: None
        \end{itemize}
      \item \textsf{void deleteRiskAssessment(string id)}
        \begin{itemize}[noitemsep,nolistsep]
           \item Effect: Delete a \ref{data:Datatypesorghl7fhirr5modelRiskAssessment} instance with the specified id. 
           \item Sequence Diagrams: None
        \end{itemize}
      \item \textsf{\ref{data:Datatypesorghl7fhirr5modelObservation} getObservation(string id)}
        \begin{itemize}[noitemsep,nolistsep]
           \item Effect: Retrieve an observation instance 
           \item Returns: The \ref{data:Datatypesorghl7fhirr5modelObservation} with the specified identifier.
           \item Sequence Diagrams: None
        \end{itemize}
      \item \textsf{\ref{data:Datatypesorghl7fhirr5modelPatient} getPatient(string id)}
        \begin{itemize}[noitemsep,nolistsep]
           \item Effect: Retrieve a patient record with the specified identifier. 
           \item Returns: The \ref{data:Datatypesorghl7fhirr5modelPatient} with the specified identifier.
           \item Sequence Diagrams: None
        \end{itemize}
      \item \textsf{\ref{data:Datatypesorghl7fhirr5modelRiskAssessment} getRiskAssessment(string id)}
        \begin{itemize}[noitemsep,nolistsep]
           \item Effect: Retrieve a \ref{data:Datatypesorghl7fhirr5modelRiskAssessment} instance with the specified id. 
           \item Returns: The \ref{data:Datatypesorghl7fhirr5modelRiskAssessment} with the specified identifier.
           \item Sequence Diagrams: None
        \end{itemize}
      \item \textsf{\ref{data:Datatypesorghl7fhirr5modelObservation} saveObservation(\ref{data:Datatypesorghl7fhirr5modelObservation} observation)}
        \begin{itemize}[noitemsep,nolistsep]
           \item Effect: Create a new \ref{data:Datatypesorghl7fhirr5modelObservation} instance or in case an existing \ref{data:Datatypesorghl7fhirr5modelObservation} instance exists with the same identifier, update this \ref{data:Datatypesorghl7fhirr5modelObservation} instance. 
           \item Returns: The \ref{data:Datatypesorghl7fhirr5modelObservation} object.
           \item Sequence Diagrams: \cref{diag:Interaction:Incomingsensordata}
        \end{itemize}
      \item \textsf{\ref{data:Datatypesorghl7fhirr5modelPatient} savePatient(\ref{data:Datatypesorghl7fhirr5modelPatient} patient)}
        \begin{itemize}[noitemsep,nolistsep]
           \item Effect: Create a new \ref{data:Datatypesorghl7fhirr5modelPatient} instance or update an existing \ref{data:Datatypesorghl7fhirr5modelPatient} instance. 
           \item Returns: The updated \ref{data:Datatypesorghl7fhirr5modelPatient} object.
           \item Sequence Diagrams: None
        \end{itemize}
      \item \textsf{\ref{data:Datatypesorghl7fhirr5modelRiskAssessment} saveRiskAssessment(\ref{data:Datatypesorghl7fhirr5modelRiskAssessment} riskAssessment)}
        \begin{itemize}[noitemsep,nolistsep]
           \item Effect: Create a new \ref{data:Datatypesorghl7fhirr5modelRiskAssessment} instance or update an existing \ref{data:Datatypesorghl7fhirr5modelRiskAssessment} instance. 
           \item Returns: The updated \ref{data:Datatypesorghl7fhirr5modelRiskAssessment} object.
           \item Sequence Diagrams: None
        \end{itemize}
      \item \textsf{List\textless{}\ref{data:Datatypesorghl7fhirr5modelObservation}\textgreater{} searchObservation(\ref{data:DatatypesQuery}\textless{}\ref{data:DatatypesDate}\textgreater{} date, \ref{data:DatatypesQuery}\textless{}String\textgreater{} dataAbsentReason, \ref{data:DatatypesQuery}\textless{}\ref{data:Datatypesorghl7fhirr5modelPatient}\textgreater{} subject, \ref{data:DatatypesQuery}\textless{}String\textgreater{} valueConcept, \ref{data:DatatypesQuery}\textless{}\ref{data:DatatypesDate}\textgreater{} valueDate, \ref{data:DatatypesQuery}\textless{}\ref{data:Datatypesorghl7fhirr5modelObservation}\textgreater{} derivedFrom, \ref{data:DatatypesQuery}\textless{}\ref{data:Datatypesorghl7fhirr5modelPatient}\textgreater{} patient, \ref{data:DatatypesQuery}\textless{}\ref{data:Datatypesorghl7fhirr5modelQuantity}\textgreater{} valueQuantity, \ref{data:DatatypesQuery}\textless{}String\textgreater{} identifier, \ref{data:DatatypesQuery}\textless{}String\textgreater{} performer, \ref{data:DatatypesQuery}\textless{}String\textgreater{} method, \ref{data:DatatypesQuery}\textless{}String\textgreater{} category, \ref{data:DatatypesQuery}\textless{}String\textgreater{} device, \ref{data:DatatypesQuery}\textless{}\ref{data:Datatypesorghl7fhirr5modelObservationStatus}\textgreater{} status)}
        \begin{itemize}[noitemsep,nolistsep]
           \item Effect: Search for \ref{data:Datatypesorghl7fhirr5modelObservation} instances matching the specified query criteria. 
           \item Returns: The list of \ref{data:Datatypesorghl7fhirr5modelObservation}s matching the specified search criteria.
           \item Sequence Diagrams: None
        \end{itemize}
      \item \textsf{List\textless{}\ref{data:Datatypesorghl7fhirr5modelPatient}\textgreater{} searchPatient(\ref{data:DatatypesQuery}\textless{}\ref{data:DatatypesDate}\textgreater{} birthDate, \ref{data:DatatypesQuery}\textless{}Boolean\textgreater{} deceased, \ref{data:DatatypesQuery}\textless{}String\textgreater{} addressState, \ref{data:DatatypesQuery}\textless{}String\textgreater{} administrativeGender, \ref{data:DatatypesQuery}\textless{}String\textgreater{} link, \ref{data:DatatypesQuery}\textless{}String\textgreater{} language, \ref{data:DatatypesQuery}\textless{}String\textgreater{} addressCountry, \ref{data:DatatypesQuery}\textless{}\ref{data:DatatypesDate}\textgreater{} deathDate, \ref{data:DatatypesQuery}\textless{}String\textgreater{} phonetic, \ref{data:DatatypesQuery}\textless{}String\textgreater{} telecom, \ref{data:DatatypesQuery}\textless{}String\textgreater{} addressCity, \ref{data:DatatypesQuery}\textless{}String\textgreater{} email, \ref{data:DatatypesQuery}\textless{}String\textgreater{} given, \ref{data:DatatypesQuery}\textless{}String\textgreater{} identifier, \ref{data:DatatypesQuery}\textless{}String\textgreater{} address, \ref{data:DatatypesQuery}\textless{}String\textgreater{} generalPractitioner, \ref{data:DatatypesQuery}\textless{}Boolean\textgreater{} active, \ref{data:DatatypesQuery}\textless{}String\textgreater{} addressPostalCode, \ref{data:DatatypesQuery}\textless{}String\textgreater{} phone, \ref{data:DatatypesQuery}\textless{}String\textgreater{} organizationCustodian, \ref{data:DatatypesQuery}\textless{}String\textgreater{} name, \ref{data:DatatypesQuery}\textless{}String\textgreater{} family)}
        \begin{itemize}[noitemsep,nolistsep]
           \item Effect: Search for \ref{data:Datatypesorghl7fhirr5modelPatient} instances matching the specified query criteria. 
           \item Returns: The list of \ref{data:Datatypesorghl7fhirr5modelPatient}s matching the specified search criteria.
           \item Sequence Diagrams: None
        \end{itemize}
      \item \textsf{List\textless{}\ref{data:Datatypesorghl7fhirr5modelRiskAssessment}\textgreater{} searchRiskAssessment(\ref{data:DatatypesQuery}\textless{}\ref{data:DatatypesDate}\textgreater{} date, \ref{data:DatatypesQuery}\textless{}String\textgreater{} identifier, \ref{data:DatatypesQuery}\textless{}String\textgreater{} performer, \ref{data:DatatypesQuery}\textless{}String\textgreater{} method, \ref{data:DatatypesQuery}\textless{}\ref{data:Datatypesorghl7fhirr5modelRange}\textgreater{} probability, \ref{data:DatatypesQuery}\textless{}\ref{data:Datatypesorghl7fhirr5modelPatient}\textgreater{} subject, \ref{data:DatatypesQuery}\textless{}String\textgreater{} condition, \ref{data:DatatypesQuery}\textless{}\ref{data:Datatypesorghl7fhirr5modelPatient}\textgreater{} patient, \ref{data:DatatypesQuery}\textless{}\ref{data:DatatypesBigDecimal}\textgreater{} risk)}
        \begin{itemize}[noitemsep,nolistsep]
           \item Effect: Search for a \ref{data:Datatypesorghl7fhirr5modelRiskAssessment} instance matching the provided query criteria. 
           \item Returns: The list of \ref{data:Datatypesorghl7fhirr5modelRiskAssessment}s matching the specified search criteria.
           \item Sequence Diagrams: None
        \end{itemize}
    \end{itemize}
      \item[Diagrams:] None
    \end{description}

  %%%%%%%% JobMgmt
  \subsection{JobMgmt}\label{int:InterfacesJobMgmt}
    \begin{description}[noitemsep,nolistsep]
      \item[Provided by:] \iconcomponent{}~\vpett{\nameref{comp:ComponentseHealthPlatformRiskEstimationCombiner}}
      \item[Required by:] \iconcomponent{}~\vpett{\nameref{comp:ComponentseHealthPlatformRiskEstimationScheduler}} 
      \item[Operations:] ~
    \begin{itemize}[noitemsep,nolistsep,leftmargin=-.25cm]
      \item \textsf{void addJobSet(\ref{data:DatatypesRiskEstimationId} id, List\textless{}Tuple\textless{}\ref{data:DatatypesClinicalModelJobId}, Double\textgreater{}\textgreater{} jobIds)}
        \begin{itemize}[noitemsep,nolistsep]
           \item Effect: Adds a set of identifiers for jobs belonging to a
single risk estimation identified by id. The partial
results of each clinical model computation job identified by an
element in jobIds have to be combined in order to find
the final result of the risk estimation as a whole. 
           \item Sequence Diagrams: \cref{diag:Interaction:Riskestimationprocess}
        \end{itemize}
    \end{itemize}
      \item[Diagrams:] None
    \end{description}

  %%%%%%%% KVStore
  \subsection{KVStore}\label{int:InterfacesKVStore}
    \begin{description}[noitemsep,nolistsep]
      \item[Provided by:] \iconcomponent{}~\vpett{\nameref{comp:ComponentsKVStorageService}}
      \item[Required by:] \iconcomponent{}~\vpett{\nameref{mod:ComponentseHealthPlatformRiskEstimationProcessorDataRiskAssessment}}, \iconcomponent{}~\vpett{\nameref{comp:ComponentseHealthPlatformRiskEstimationProcessor}}, \iconcomponent{}~\vpett{\nameref{comp:ComponentseHealthPlatform}} 
      \item[Operations:] ~
    \begin{itemize}[noitemsep,nolistsep,leftmargin=-.25cm]
      \item \textsf{\ref{data:DatatypesObject} retrieve(String key) throws \ref{ex:ExceptionsIOException}}
        \begin{itemize}[noitemsep,nolistsep]
           \item Effect: Retrieve the object with the provided key from the Key-Value storage. 
           \item Returns: The object retrieved from storage.
           \item Sequence Diagrams: \cref{diag:Interaction:Computeclinicalmodelresult}
        \end{itemize}
      \item \textsf{void store(String key, \ref{data:DatatypesObject} value) throws \ref{ex:ExceptionsIOException}}
        \begin{itemize}[noitemsep,nolistsep]
           \item Effect: Store the provided object in the Key-Value storage under the provided key. 
           \item Sequence Diagrams: \cref{diag:Interaction:Computeclinicalmodelresult}
        \end{itemize}
    \end{itemize}
      \item[Diagrams:] None
    \end{description}

  %%%%%%%% LaunchRiskEstimation
  \subsection{LaunchRiskEstimation}\label{int:InterfacesLaunchRiskEstimation}
    \begin{description}[noitemsep,nolistsep]
      \item[Provided by:] \iconcomponent{}~\vpett{\nameref{comp:ComponentseHealthPlatformRiskEstimationScheduler}}, \iconcomponent{}~\vpett{\nameref{comp:ComponentseHealthPlatform}}
      \item[Required by:] \iconcomponent{}~\vpett{\nameref{comp:ComponentsPatientGatewayEndpoint}} 
      \item[Operations:] ~
    \begin{itemize}[noitemsep,nolistsep,leftmargin=-.25cm]
      \item \textsf{void launchRiskEstimation(\ref{data:DatatypesPatientId} patientId, \ref{data:DatatypesSensorDataPackage} newSensorData, \ref{data:DatatypesTimestamp} receivedAt)}
        \begin{itemize}[noitemsep,nolistsep]
           \item Effect: The \vpett{\nameref{comp:ComponentseHealthPlatformRiskEstimationScheduler}} will fetch the clinical models and their configurations associated to the patient identified by patientId from storage using the \vpett{\nameref{comp:ComponentseHealthPlatformClinicalModelCache}}, notify the \vpett{\nameref{comp:ComponentseHealthPlatformRiskEstimationCombiner}} of the different jobs that will be performed for a single risk estimation and schedule the individual jobs in its queue.
In normal modus, queued jobs are returned in FIFO order. In overload modus, the system switches to dynamic priority: earliest deadline first and enqueues jobs of patient with a high risk level with an earlier deadline (2 instead of 5 minutes) to prioritize them over patients with lower risk levels.
The \ref{data:DatatypesSensorDataPackage} newSensorData is passed because its arrival triggered the risk estimation. A risk level is estimated based on the computation of these clinical models. For the computation of the clinical models, the given sensor data newSensorData (which is the new sensor data that was received) is used and other required data is fetched  from the respective databases if needed. The given time-stamp receivedAt is used in order to avoid fetching the new sensor data from the database. This time-stamp represents the time at which the new sensor data was received. 
           \item Sequence Diagrams: \cref{diag:Interaction:Riskestimationprocess,diag:Interaction:Incomingsensordata}
        \end{itemize}
    \end{itemize}
      \item[Diagrams:] None
    \end{description}

  %%%%%%%% MLaaSAPI
  \subsection{MLaaSAPI}\label{int:InterfacesMLaaSAPI}
    \begin{description}[noitemsep,nolistsep]
      \item[Provided by:] \iconcomponent{}~\vpett{\nameref{mod:ModulesMLaaSLibrary}}
      \item[Required by:] \iconcomponent{}~\vpett{\nameref{mod:ComponentseHealthPlatformMLModelManagerMLModelUpdateProcessor}} 
           \item[Description]: This API provides a number of operations for managing the models deployed on the \vpett{\nameref{comp:ComponentsMLaaS}} cloud.
      \item[Operations:] ~
    \begin{itemize}[noitemsep,nolistsep,leftmargin=-.25cm]
      \item \textsf{\ref{data:DatatypesCredential} authenticate() throws \ref{ex:ExceptionsInvalidAPIKeyException}}
        \begin{itemize}[noitemsep,nolistsep]
           \item Effect: Construct a \ref{data:DatatypesCredential} object for authenticating requests with the \vpett{\nameref{comp:ComponentsMLaaS}} service provider. This operation uses the previously configured API key to construct the \ref{data:DatatypesCredential} object. If such an API key is not configured, an exception will be thrown. 
           \item Returns: Authentication token for subsequent requests.
           \item Sequence Diagrams: \cref{diag:Interaction:MLModelSynchronization,diag:Interaction:ProcessingMLModelupdates}
        \end{itemize}
      \item \textsf{void delete(\ref{data:DatatypesMLModel} model, \ref{data:DatatypesCredential} token) throws \ref{ex:ExceptionsAuthorizationException}}
        \begin{itemize}[noitemsep,nolistsep]
           \item Effect: Delete a model from the online \vpett{\nameref{comp:ComponentsMLaaS}} service. 
           \item Sequence Diagrams: None
        \end{itemize}
      \item \textsf{\ref{data:DatatypesMLModelDescription} getModelDescription(List\textless{}\ref{data:DatatypesMLModelDescription}\textgreater{} descriptions, String modelId)}
        \begin{itemize}[noitemsep,nolistsep]
           \item Effect: This method retrieves the correct \ref{data:DatatypesMLModelDescription} that corresponds with the \ref{data:DatatypesMLModel} identified by the provided modelId. 
           \item Returns: The \ref{data:DatatypesMLModelDescription} corresponding with the \ref{data:DatatypesMLModel} identified by the modelId.
           \item Sequence Diagrams: None
        \end{itemize}
      \item \textsf{String getModelId(\ref{data:DatatypesMLModelDescription} description)}
        \begin{itemize}[noitemsep,nolistsep]
           \item Effect: Extract the modelId from the \ref{data:DatatypesMLModelDescription}. 
           \item Returns: String with the modelId of the \ref{data:DatatypesMLModel} specified in the \ref{data:DatatypesMLModelDescription}.
           \item Sequence Diagrams: \cref{diag:Interaction:MLModelSynchronization}
        \end{itemize}
      \item \textsf{void register(\ref{data:DatatypesMLModel} model, \ref{data:DatatypesCredential} token) throws \ref{ex:ExceptionsAuthorizationException}}
        \begin{itemize}[noitemsep,nolistsep]
           \item Effect: Register a new model with the \vpett{\nameref{comp:ComponentsMLaaS}} service. 
           \item Sequence Diagrams: None
        \end{itemize}
      \item \textsf{void setAPIKey(String apiKey) throws \ref{ex:ExceptionsAuthenticationException}}
        \begin{itemize}[noitemsep,nolistsep]
           \item Effect: Set the credentials for authenticating with the ML-as-a-Service provider. 
           \item Sequence Diagrams: None
        \end{itemize}
    \end{itemize}
      \item[Diagrams:] None
    \end{description}

  %%%%%%%% MLaaSService
  \subsection{MLaaSService}\label{int:InterfacesMLaaSService}
    \begin{description}[noitemsep,nolistsep]
      \item[Provided by:] \iconcomponent{}~\vpett{\nameref{comp:ComponentsMLaaS}}
      \item[Required by:] \iconcomponent{}~\vpett{\nameref{comp:ComponentseHealthPlatformMLModelManager}}, \iconcomponent{}~\vpett{\nameref{mod:ComponentseHealthPlatformMLModelManagerMLModelUpdateProcessor}}, \iconcomponent{}~\vpett{\nameref{comp:ComponentseHealthPlatform}} 
           \item[Description]: This interface offiers operations for interacting with \vpett{\nameref{comp:ComponentsMLaaS}} service providers.
      \item[Operations:] ~
    \begin{itemize}[noitemsep,nolistsep,leftmargin=-.25cm]
      \item \textsf{List\textless{}\ref{data:DatatypesMLModelUpdate}\textgreater{} fetchModelUpdates(\ref{data:DatatypesCredential} credential, \ref{data:DatatypesMLModelDescription} model)}
        \begin{itemize}[noitemsep,nolistsep]
           \item Effect: Fetch a list of model updates to apply to the local \ref{data:DatatypesMLModel}. 
           \item Returns: The list of MLModelUpdates to apply locally.
           \item Sequence Diagrams: \cref{diag:Interaction:ProcessingMLModelupdates}
        \end{itemize}
      \item \textsf{List\textless{}\ref{data:DatatypesMLModelDescription}\textgreater{} getAvailableModels(\ref{data:DatatypesCredential} credential)}
        \begin{itemize}[noitemsep,nolistsep]
           \item Effect: Get the list of available ML models to use for making predictions on incoming sensor data. 
           \item Returns: The list of available MLModels.
           \item Sequence Diagrams: \cref{diag:Interaction:MLModelSynchronization}
        \end{itemize}
      \item \textsf{void pushModelUpdates(\ref{data:DatatypesCredential} credential, \ref{data:DatatypesMLModelDescription} model, List\textless{}\ref{data:DatatypesMLModelUpdate}\textgreater{} updates)}
        \begin{itemize}[noitemsep,nolistsep]
           \item Effect: Send a list of MLModelUpdates to the online service. 
           \item Sequence Diagrams: None
        \end{itemize}
      \item \textsf{\ref{data:DatatypesMLModel} retrieveModel(\ref{data:DatatypesCredential} credential, \ref{data:DatatypesMLModelDescription} model)}
        \begin{itemize}[noitemsep,nolistsep]
           \item Effect: Retrieve a specific \ref{data:DatatypesMLModel} corresponding with the provided \ref{data:DatatypesMLModelDescription}. 
           \item Returns: The requested \ref{data:DatatypesMLModel}.
           \item Sequence Diagrams: \cref{diag:Interaction:MLModelSynchronization}
        \end{itemize}
    \end{itemize}
      \item[Diagrams:] None
    \end{description}

  %%%%%%%% MLAPI
  \subsection{MLAPI}\label{int:InterfacesMLAPI}
    \begin{description}[noitemsep,nolistsep]
      \item[Provided by:] \iconcomponent{}~\vpett{\nameref{mod:ModulesMLLibrary}}
      \item[Required by:] \iconcomponent{}~\vpett{\nameref{mod:ComponentseHealthPlatformRiskEstimationProcessorDataRiskAssessment}}, \iconcomponent{}~\vpett{\nameref{mod:ComponentseHealthPlatformMLModelManagerModelUpdater}} 
           \item[Description]: This is the API provided by the \vpett{\nameref{mod:ModulesMLLibrary}} to use \ref{data:DatatypesMLModel}s locally for making predictions and processing updates to these models.
      \item[Operations:] ~
    \begin{itemize}[noitemsep,nolistsep,leftmargin=-.25cm]
      \item \textsf{\ref{data:DatatypesMLModel} applyUpdate(\ref{data:DatatypesMLModel} model, \ref{data:DatatypesMLModelUpdate} update)}
        \begin{itemize}[noitemsep,nolistsep]
           \item Effect: Applies the provided \ref{data:DatatypesMLModelUpdate} on the provided \ref{data:DatatypesMLModel} and returns the resulting \ref{data:DatatypesMLModel}. 
           \item Returns: The resulting \ref{data:DatatypesMLModel} after applying the requested \ref{data:DatatypesMLModelUpdate}.
           \item Sequence Diagrams: \cref{diag:Interaction:ProcessingMLModelupdates}
        \end{itemize}
      \item \textsf{\ref{data:DatatypesPrediction} predict(\ref{data:DatatypesSensorDataPackage} sensorData)}
        \begin{itemize}[noitemsep,nolistsep]
           \item Effect: Make a prediction with the active \ref{data:DatatypesMLModel} on the provided sensordata. 
           \item Returns: The prediction based on the provided sensor data.
           \item Sequence Diagrams: \cref{diag:Interaction:Computeclinicalmodelresult}
        \end{itemize}
      \item \textsf{void verify(\ref{data:DatatypesMLModel} model) throws \ref{ex:ExceptionsInvalidMLModelException}}
        \begin{itemize}[noitemsep,nolistsep]
           \item Effect: Verifies that a provided \ref{data:DatatypesMLModel} is working correctly after applying \ref{data:DatatypesMLModelUpdate}s. 
           \item Sequence Diagrams: \cref{diag:Interaction:ProcessingMLModelupdates}
        \end{itemize}
    \end{itemize}
      \item[Diagrams:] None
    \end{description}

  %%%%%%%% MLModelMgmt
  \subsection{MLModelMgmt}\label{int:InterfacesMLModelMgmt}
    \begin{description}[noitemsep,nolistsep]
      \item[Provided by:] \iconcomponent{}~\vpett{\nameref{comp:ComponentseHealthPlatformMLModelManager}}, \iconcomponent{}~\vpett{\nameref{mod:ComponentseHealthPlatformMLModelManagerMLModelStorageManager}}
      \item[Required by:] \iconcomponent{}~\vpett{\nameref{mod:ComponentseHealthPlatformRiskEstimationProcessorMLModelRetrieval}}, \iconcomponent{}~\vpett{\nameref{comp:ComponentseHealthPlatformRiskEstimationProcessor}} 
           \item[Description]: The interface provides operations for managing different \ref{data:DatatypesMLModel}s locally, querying available models for making predictions on newly incoming sensor data.
      \item[Operations:] ~
    \begin{itemize}[noitemsep,nolistsep,leftmargin=-.25cm]
      \item \textsf{\ref{data:DatatypesMLModel} fetchModel(String modelId) throws \ref{ex:ExceptionsNoSuchMLModelException}}
        \begin{itemize}[noitemsep,nolistsep]
           \item Effect: Fetch the \ref{data:DatatypesMLModel} corresponding with the provided modelId. 
           \item Returns: The requested \ref{data:DatatypesMLModel}.
           \item Sequence Diagrams: \cref{diag:Interaction:Computeclinicalmodelresult}
        \end{itemize}
      \item \textsf{List\textless{}String\textgreater{} getAvailableModels()}
        \begin{itemize}[noitemsep,nolistsep]
           \item Effect: Get the list of available \ref{data:DatatypesMLModel}s that can be used for making \ref{data:DatatypesPrediction}s locally. 
           \item Returns: The list of identifiers of the available \ref{data:DatatypesMLModel}s that can be used.
           \item Sequence Diagrams: None
        \end{itemize}
      \item \textsf{void synchronizeModels()}
        \begin{itemize}[noitemsep,nolistsep]
           \item Effect: Synchronize the list of locally-stored \ref{data:DatatypesMLModel}s by checking for newly-available \ref{data:DatatypesMLModel}s that are not stored in the local \vpett{\nameref{mod:ComponentseHealthPlatformMLModelManagerMLModelRepository}} and retrieving these \ref{data:DatatypesMLModel}s from the \vpett{\nameref{comp:ComponentsMLaaS}} provider. 
           \item Sequence Diagrams: \cref{diag:Interaction:MLModelSynchronization}
        \end{itemize}
      \item \textsf{void updateModel(String modelId) throws \ref{ex:ExceptionsNoSuchMLModelException}}
        \begin{itemize}[noitemsep,nolistsep]
           \item Effect: Update the model with the specified id by fetching \ref{data:DatatypesMLModelUpdate}s and applying them to the model. 
           \item Sequence Diagrams: \cref{diag:Interaction:ProcessingMLModelupdates}
        \end{itemize}
    \end{itemize}
      \item[Diagrams:] None
    \end{description}

  %%%%%%%% MLModelStorage
  \subsection{MLModelStorage}\label{int:InterfacesMLModelStorage}
    \begin{description}[noitemsep,nolistsep]
      \item[Provided by:] \iconcomponent{}~\vpett{\nameref{mod:ComponentseHealthPlatformMLModelManagerMLModelRepository}}
      \item[Required by:] \iconcomponent{}~\vpett{\nameref{mod:ComponentseHealthPlatformMLModelManagerMLModelStorageManager}}, \iconcomponent{}~\vpett{\nameref{mod:ComponentseHealthPlatformMLModelManagerMLModelUpdateProcessor}} 
      \item[Operations:] ~
    \begin{itemize}[noitemsep,nolistsep,leftmargin=-.25cm]
      \item \textsf{boolean isAvailable(String modelId)}
        \begin{itemize}[noitemsep,nolistsep]
           \item Effect: Check whether an \ref{data:DatatypesMLModel} (identified by the modelId) is available in the local \vpett{\nameref{mod:ComponentseHealthPlatformMLModelManagerMLModelRepository}}. 
           \item Returns: Boolean indicating if the specified model is present in the repository.
           \item Sequence Diagrams: \cref{diag:Interaction:MLModelSynchronization,diag:Interaction:ProcessingMLModelupdates}
        \end{itemize}
      \item \textsf{\ref{data:DatatypesMLModel} retrieveModel(String modelId)}
        \begin{itemize}[noitemsep,nolistsep]
           \item Effect: Retrieve the \ref{data:DatatypesMLModel} with the specified id. 
           \item Returns: The requested \ref{data:DatatypesMLModel}.
           \item Sequence Diagrams: \cref{diag:Interaction:ProcessingMLModelupdates}
        \end{itemize}
      \item \textsf{void storeModel(\ref{data:DatatypesMLModel} model)}
        \begin{itemize}[noitemsep,nolistsep]
           \item Effect: Store an \ref{data:DatatypesMLModel} in the local \vpett{\nameref{mod:ComponentseHealthPlatformMLModelManagerMLModelRepository}} 
           \item Sequence Diagrams: \cref{diag:Interaction:MLModelSynchronization,diag:Interaction:ProcessingMLModelupdates}
        \end{itemize}
    \end{itemize}
      \item[Diagrams:] None
    \end{description}

  %%%%%%%% MLModelSync
  \subsection{MLModelSync}\label{int:InterfacesMLModelSync}
    \begin{description}[noitemsep,nolistsep]
      \item[Provided by:] \iconcomponent{}~\vpett{\nameref{mod:ComponentseHealthPlatformMLModelManagerMLModelUpdateProcessor}}
      \item[Required by:] \iconcomponent{}~\vpett{\nameref{mod:ComponentseHealthPlatformMLModelManagerMLModelStorageManager}} 
      \item[Operations:] ~
    \begin{itemize}[noitemsep,nolistsep,leftmargin=-.25cm]
      \item \textsf{void fetchNewMLModels()}
        \begin{itemize}[noitemsep,nolistsep]
           \item Effect: Fetches newly available \ref{data:DatatypesMLModel}s and store them in the \vpett{\nameref{mod:ComponentseHealthPlatformMLModelManagerMLModelRepository}}. 
           \item Sequence Diagrams: \cref{diag:Interaction:MLModelSynchronization}
        \end{itemize}
      \item \textsf{void processMLModelUpdates(String modelId) throws \ref{ex:ExceptionsNoSuchMLModelException}}
        \begin{itemize}[noitemsep,nolistsep]
           \item Effect: Retrieve and process the updates for the specified \ref{data:DatatypesMLModel} and store the updated \ref{data:DatatypesMLModel} in the \vpett{\nameref{mod:ComponentseHealthPlatformMLModelManagerMLModelRepository}}. 
           \item Sequence Diagrams: \cref{diag:Interaction:ProcessingMLModelupdates}
        \end{itemize}
    \end{itemize}
      \item[Diagrams:] None
    \end{description}

  %%%%%%%% MLModelUpdateMgmt
  \subsection{MLModelUpdateMgmt}\label{int:InterfacesMLModelUpdateMgmt}
    \begin{description}[noitemsep,nolistsep]
      \item[Provided by:] \iconcomponent{}~\vpett{\nameref{mod:ComponentseHealthPlatformMLModelManagerModelUpdater}}
      \item[Required by:] \iconcomponent{}~\vpett{\nameref{mod:ComponentseHealthPlatformMLModelManagerMLModelUpdateProcessor}} 
      \item[Operations:] ~
    \begin{itemize}[noitemsep,nolistsep,leftmargin=-.25cm]
      \item \textsf{\ref{data:DatatypesMLModel} processMLModelUpdates(\ref{data:DatatypesMLModel} model, List\textless{}\ref{data:DatatypesMLModelUpdate}\textgreater{} updates)}
        \begin{itemize}[noitemsep,nolistsep]
           \item Effect: Requests the \vpett{\nameref{mod:ComponentseHealthPlatformMLModelManagerModelUpdater}} to apply the list of  \ref{data:DatatypesMLModelUpdate}s on the provided \ref{data:DatatypesMLModel}. 
           \item Returns: The \ref{data:DatatypesMLModel} resulting from the application of the list of \ref{data:DatatypesMLModelUpdate}s on the initial \ref{data:DatatypesMLModel}.
           \item Sequence Diagrams: \cref{diag:Interaction:ProcessingMLModelupdates}
        \end{itemize}
    \end{itemize}
      \item[Diagrams:] None
    \end{description}

  %%%%%%%% OS Mgmt
  \subsection{OS Mgmt}\label{int:InterfacesOSMgmt}
    \begin{description}[noitemsep,nolistsep]
      \item[Provided by:] \iconcomponent{}~\vpett{\nameref{comp:ComponentsMobileOS}}
      \item[Required by:] \iconcomponent{}~\vpett{\nameref{mod:ComponentsPatientGatewayAppAppEndpoint}}, \iconcomponent{}~\vpett{\nameref{comp:ComponentsPatientGatewayApp}} 
           \item[Description]: Interface for mobile OS of users phone and the app. 
      \item[Operations:] ~
    \begin{itemize}[noitemsep,nolistsep,leftmargin=-.25cm]
      \item \textsf{void sendBatteryNotif()}
        \begin{itemize}[noitemsep,nolistsep]
           \item Effect: send a notification that the battery of the smartphone is low 
           \item Sequence Diagrams: \cref{diag:Interaction:Patientgatewayfails}
        \end{itemize}
    \end{itemize}
      \item[Diagrams:] None
    \end{description}

  %%%%%%%% OtherDataMgmt
  \subsection{OtherDataMgmt}\label{int:InterfacesOtherDataMgmt}
    \begin{description}[noitemsep,nolistsep]
      \item[Provided by:] \iconcomponent{}~\vpett{\nameref{comp:ComponentseHealthPlatformPatientDataStorage}}
      \item[Required by:] \iconcomponent{}~\vpett{\nameref{comp:ComponentseHealthPlatformAlertManager}}, \iconcomponent{}~\vpett{\nameref{comp:ComponentseHealthPlatformClinicalJobCreator}}, \iconcomponent{}~\vpett{\nameref{comp:ComponentseHealthPlatformClinicalModelCache}}, \iconcomponent{}~\vpett{\nameref{comp:ComponentseHealthPlatformRiskEstimationCombiner}}, \iconcomponent{}~\vpett{\nameref{comp:ComponentseHealthPlatformRiskEstimationScheduler}} 
      \item[Operations:] ~
    \begin{itemize}[noitemsep,nolistsep,leftmargin=-.25cm]
      \item \textsf{\ref{data:DatatypesPatientStatus} getPatientStatus(\ref{data:DatatypesPatientId} patientId) throws \ref{ex:ExceptionsNoSuchPatientException}}
        \begin{itemize}[noitemsep,nolistsep]
           \item Effect: Fetch and return the status of the patient identified by the patientId. 
           \item Returns: The status of the patient.
           \item Sequence Diagrams: \cref{diag:Interaction:av2CommunicationchannelbetweenPMSandPGgoesdownDetectionbytimeoutBackend}
        \end{itemize}
      \item \textsf{void setEstimatedPatientStatus(\ref{data:DatatypesPatientId} patientId, \ref{data:DatatypesPatientStatus} estimatedStatus, \ref{data:DatatypesTimestamp} estimationTime) throws \ref{ex:ExceptionsNoSuchPatientException}}
        \begin{itemize}[noitemsep,nolistsep]
           \item Effect: Update the patient status estimation of the patient identified by patientId to the given value estimatedStatus and update the time of estimation to estimationTime. The time of estimation is the time at which the corresponding estimation job for this patient was completed. If patient's estimated risk changed the appropriate parties are notified. 
           \item Sequence Diagrams: \cref{diag:Interaction:Incomingsensordata}
        \end{itemize}
    \end{itemize}
      \item[Diagrams:] None
    \end{description}

  %%%%%%%% PatientDataMgmt
  \subsection{PatientDataMgmt}\label{int:InterfacesPatientDataMgmt}
    \begin{description}[noitemsep,nolistsep]
      \item[Provided by:] \iconcomponent{}~\vpett{\nameref{comp:ComponentsPatientGatewayEndpoint}}
      \item[Required by:] \iconcomponent{}~\vpett{\nameref{mod:ComponentsPatientGatewayAppAppEndpoint}}, \iconcomponent{}~\vpett{\nameref{comp:ComponentsPatientGatewayApp}} 
      \item[Operations:] ~
    \begin{itemize}[noitemsep,nolistsep,leftmargin=-.25cm]
      \item \textsf{boolean addTPToken(List\textless{}\ref{data:DatatypesTokenPackage}\textgreater{} newTokens)}
        \begin{itemize}[noitemsep,nolistsep]
           \item Effect: A function for storing tokens to the pms system 
           \item Returns: boolean wheter the call succeeded
           \item Sequence Diagrams: \cref{diag:Interaction:Storetokens}
        \end{itemize}
      \item \textsf{boolean flushBuffer(\ref{data:DatatypesSensorDataPackage} sensordata)}
        \begin{itemize}[noitemsep,nolistsep]
           \item Effect: Flush data in the buffer of the app to the PMS. 
           \item Returns: Returns true if data from the buffer gets delivered succesfully, false otherwise.
           \item Sequence Diagrams: \cref{diag:Interaction:Exitdegradedmode}
        \end{itemize}
      \item \textsf{boolean ping()}
        \begin{itemize}[noitemsep,nolistsep]
           \item Effect: ping to see if connection is okay
 
           \item Returns: connection succefull
           \item Sequence Diagrams: \cref{diag:Interaction:Ping}
        \end{itemize}
      \item \textsf{void putAlert()}
        \begin{itemize}[noitemsep,nolistsep]
           \item Effect: send an alert to the patient gateway 
           \item Sequence Diagrams: \cref{diag:Interaction:Patientgatewayfails}
        \end{itemize}
      \item \textsf{boolean putPatientData(\ref{data:DatatypesSensorDataPackage} data, \ref{data:DatatypesPatientId} patientid)}
        \begin{itemize}[noitemsep,nolistsep]
           \item Effect: Wearable unit data is sent to patient gateway endpoint from the patient gateway app. Endpoint sends an acknowlegdement back if it has received the data. 
           \item Returns: the sensor data from the patient
           \item Sequence Diagrams: \cref{diag:Interaction:Incomingsensordata,diag:Interaction:av2CommunicationchannelbetweenPMSandPGgoesdownDetectionduringCommunicationAppSide}
        \end{itemize}
    \end{itemize}
      \item[Diagrams:] None
    \end{description}

  %%%%%%%% PatientDataPush
  \subsection{PatientDataPush}\label{int:InterfacesPatientDataPush}
    \begin{description}[noitemsep,nolistsep]
      \item[Provided by:] \iconcomponent{}~\vpett{\nameref{mod:ComponentsPatientGatewayAppAppEndpoint}}, \iconcomponent{}~\vpett{\nameref{comp:ComponentsPatientGatewayApp}}
      \item[Required by:] \iconcomponent{}~\vpett{\nameref{comp:ComponentsPatientGatewayEndpoint}} 
           \item[Description]: Interface allowing the server to push data to the client device
      \item[Operations:] ~
    \begin{itemize}[noitemsep,nolistsep,leftmargin=-.25cm]
      \item \textsf{\ref{data:DatatypesSensorDataPackage} getNewSensorData()}
        \begin{itemize}[noitemsep,nolistsep]
           \item Effect: get the latest sensor data 
           \item Returns: the latest sensor data package
           \item Sequence Diagrams: None
        \end{itemize}
      \item \textsf{boolean pushNewSensorWaitTime(Map\textless{}string, int\textgreater{} sensorTimings)}
        \begin{itemize}[noitemsep,nolistsep]
           \item Effect: used to update the timings for when specific sensor data needs to be send 
           \item Returns: succes or not
           \item Sequence Diagrams: \cref{diag:Interaction:Updategatewaypushwaittime}
        \end{itemize}
    \end{itemize}
      \item[Diagrams:] None
    \end{description}

  %%%%%%%% PatientRecordMgmt
  \subsection{PatientRecordMgmt}\label{int:InterfacesPatientRecordMgmt}
    \begin{description}[noitemsep,nolistsep]
      \item[Provided by:] \iconcomponent{}~\vpett{\nameref{comp:ComponentseHealthPlatformHISManager}}
      \item[Required by:] \iconcomponent{}~\vpett{\nameref{comp:ComponentseHealthPlatformRiskEstimationCombiner}} 
      \item[Operations:] ~
    \begin{itemize}[noitemsep,nolistsep,leftmargin=-.25cm]
      \item \textsf{\ref{data:DatatypesPatientRecord} getPatientRecord(\ref{data:DatatypesPatientId} patientId) throws \ref{ex:ExceptionsNoSuchPatientRecordException}}
        \begin{itemize}[noitemsep,nolistsep]
           \item Effect: This will fetch the EHR record of the patient with given id from the \vpett{\nameref{comp:ComponentseHealthPlatformHIS}} and return it. If the \vpett{\nameref{comp:ComponentseHealthPlatformHIS}} is not available, an older (cached) copy of the patient record is returned if possible. 
           \item Returns: The requested patient record.
           \item Sequence Diagrams: None
        \end{itemize}
    \end{itemize}
      \item[Diagrams:] None
    \end{description}

  %%%%%%%% ResetTimer
  \subsection{ResetTimer}\label{int:InterfacesResetTimer}
    \begin{description}[noitemsep,nolistsep]
      \item[Provided by:] \iconcomponent{}~\vpett{\nameref{comp:ComponentseHealthPlatformCronjob}}, \iconcomponent{}~\vpett{\nameref{comp:ComponentseHealthPlatform}}
      \item[Required by:] \iconcomponent{}~\vpett{\nameref{comp:ComponentsPatientGatewayEndpoint}} 
      \item[Operations:] ~
    \begin{itemize}[noitemsep,nolistsep,leftmargin=-.25cm]
      \item \textsf{boolean resetPatientTimer(\ref{data:DatatypesPatientId} patientId, \ref{data:DatatypesTimestamp} timer)}
        \begin{itemize}[noitemsep,nolistsep]
           \item Effect: Update the timer per patient for when the patient gateway endpoint should recieve sensordata 
           \item Returns: a boolean that tells if the update was succesfull
           \item Sequence Diagrams: \cref{diag:Interaction:Incomingsensordata}
        \end{itemize}
    \end{itemize}
      \item[Diagrams:] None
    \end{description}

  %%%%%%%% Results
  \subsection{Results}\label{int:InterfacesResults}
    \begin{description}[noitemsep,nolistsep]
      \item[Provided by:] \iconcomponent{}~\vpett{\nameref{comp:ComponentseHealthPlatformRiskEstimationCombiner}}
      \item[Required by:] \iconcomponent{}~\vpett{\nameref{mod:ComponentseHealthPlatformRiskEstimationProcessorJobProcessor}}, \iconcomponent{}~\vpett{\nameref{comp:ComponentseHealthPlatformRiskEstimationProcessor}} 
      \item[Operations:] ~
    \begin{itemize}[noitemsep,nolistsep,leftmargin=-.25cm]
      \item \textsf{void addClinicalModelJobResults(\ref{data:DatatypesClinicalModelJobId} jobId, \ref{data:DatatypesClinicalModelJobResult} results)}
        \begin{itemize}[noitemsep,nolistsep]
           \item Effect: Sends the result of the performed clinical model computation (identified by the jobId) to the \vpett{\nameref{comp:ComponentseHealthPlatformRiskEstimationCombiner}} for combining with the other partial results belonging to the same risk estimation. 
           \item Sequence Diagrams: \cref{diag:Interaction:Computeclinicalmodelresult,diag:Interaction:CombiningClinicalModelJobresults}
        \end{itemize}
    \end{itemize}
      \item[Diagrams:] None
    \end{description}

  %%%%%%%% SearchToken
  \subsection{SearchToken}\label{int:InterfacesSearchToken}
    \begin{description}[noitemsep,nolistsep]
      \item[Provided by:] \iconcomponent{}~\vpett{\nameref{mod:ComponentsThirdPartyServiceTokenStorage}}
      \item[Required by:] \iconcomponent{}~\vpett{\nameref{mod:ComponentsThirdPartyServiceThirdPartySyncService}} 
           \item[Description]: the interface to search a specific token
      \item[Operations:] ~
    \begin{itemize}[noitemsep,nolistsep,leftmargin=-.25cm]
      \item \textsf{string getToken(\ref{data:DatatypesPatientId} patientId, string provider)}
        \begin{itemize}[noitemsep,nolistsep]
           \item Effect: gets a token for a patientId and provider (third party) to be able to use for the api to that third party 
           \item Returns: the token for that patient and provider
           \item Sequence Diagrams: None
        \end{itemize}
    \end{itemize}
      \item[Diagrams:] None
    \end{description}

  %%%%%%%% SensorDataMgmt
  \subsection{SensorDataMgmt}\label{int:InterfacesSensorDataMgmt}
    \begin{description}[noitemsep,nolistsep]
      \item[Provided by:] \iconcomponent{}~\vpett{\nameref{comp:ComponentseHealthPlatformHISManager}}
      \item[Required by:] \iconcomponent{}~\vpett{\nameref{comp:ComponentseHealthPlatformClinicalJobCreator}}, \iconcomponent{}~\vpett{\nameref{comp:ComponentsPatientGatewayEndpoint}}, \iconcomponent{}~\vpett{\nameref{comp:ComponentseHealthPlatformRiskEstimationScheduler}} 
      \item[Operations:] ~
    \begin{itemize}[noitemsep,nolistsep,leftmargin=-.25cm]
      \item \textsf{void addSensorData(\ref{data:DatatypesPatientId} patientId, \ref{data:DatatypesSensorDataPackage} sensorData, \ref{data:DatatypesTimestamp} receivedAt)}
        \begin{itemize}[noitemsep,nolistsep]
           \item Effect: Store the given sensor data and meta-data. 
           \item Sequence Diagrams: \cref{diag:Interaction:Incomingsensordata}
        \end{itemize}
      \item \textsf{Map\textless{}\ref{data:DatatypesTimestamp}, \ref{data:DatatypesSensorDataPackage}\textgreater{} getAllSensorDataOfPatient(\ref{data:DatatypesPatientId} patientId)}
        \begin{itemize}[noitemsep,nolistsep]
           \item Effect: This will fetch and return all sensor data belonging to the patient identified by patientId 
           \item Returns: The sensor data and the timestamp of their arrival in the system.
           \item Sequence Diagrams: None
        \end{itemize}
      \item \textsf{Map\textless{}\ref{data:DatatypesTimestamp}, \ref{data:DatatypesSensorDataPackage}\textgreater{} getAllSensorDataOfPatientBefore(\ref{data:DatatypesPatientId} patientId, \ref{data:DatatypesTimestamp} stopTime)}
        \begin{itemize}[noitemsep,nolistsep]
           \item Effect: Fetch and return all sensor data belonging to the patient identified by patientId which was received before the specified time stopTime. 
           \item Returns: The sensor data and the timestamp of their arrival in the system.
           \item Sequence Diagrams: \cref{diag:Interaction:Riskestimationprocess}
        \end{itemize}
    \end{itemize}
      \item[Diagrams:] None
    \end{description}

  %%%%%%%% ShutDown
  \subsection{ShutDown}\label{int:InterfacesShutDown}
    \begin{description}[noitemsep,nolistsep]
      \item[Provided by:] \iconcomponent{}~\vpett{\nameref{mod:ComponentsPatientGatewayAppSyncManager}}
      \item[Required by:] \iconcomponent{}~\vpett{\nameref{mod:ComponentsPatientGatewayAppAppEndpoint}} 
      \item[Operations:] ~
    \begin{itemize}[noitemsep,nolistsep,leftmargin=-.25cm]
      \item \textsf{void InitiateDegradedMode(\ref{data:DatatypesSensorDataPackage} sensordata)}
        \begin{itemize}[noitemsep,nolistsep]
           \item Effect: Ensures any data in volatile memory is pushed to the \vpett{\nameref{mod:ComponentsPatientGatewayAppBuffer}} immediately. 
           \item Sequence Diagrams: \cref{diag:Interaction:DegradedMode}
        \end{itemize}
    \end{itemize}
      \item[Diagrams:] None
    \end{description}

  %%%%%%%% SMSinterface
  \subsection{SMSinterface}\label{int:InterfacesSMSinterface}
    \begin{description}[noitemsep,nolistsep]
      \item[Provided by:] \iconcomponent{}~\vpett{\nameref{comp:ComponentsSMSgateway}}
      \item[Required by:] \iconcomponent{}~\vpett{\nameref{comp:ComponentsPatientGatewayApp}}, \iconcomponent{}~\vpett{\nameref{mod:ComponentsPatientGatewayAppSyncManager}} 
           \item[Description]: To trigger backup alerts.
      \item[Operations:] ~
    \begin{itemize}[noitemsep,nolistsep,leftmargin=-.25cm]
      \item \textsf{void sendSMS(string parameter)}
        \begin{itemize}[noitemsep,nolistsep]
           \item Effect: Called by \vpett{\nameref{mod:ComponentsPatientGatewayAppSyncManager}} to indicate sending of SMS as alternative channel when usual channel problems occur. 
           \item Sequence Diagrams: \cref{diag:Interaction:av2CommunicationchannelbetweenPMSandPGgoesdownDetectionduringCommunicationAppSide}
        \end{itemize}
    \end{itemize}
      \item[Diagrams:] None
    \end{description}

  %%%%%%%% SyncMgmt
  \subsection{SyncMgmt}\label{int:InterfacesSyncMgmt}
    \begin{description}[noitemsep,nolistsep]
      \item[Provided by:] \iconcomponent{}~\vpett{\nameref{mod:ComponentsPatientGatewayAppAppEndpoint}}
      \item[Required by:] \iconcomponent{}~\vpett{\nameref{mod:ComponentsPatientGatewayAppSyncManager}} 
           \item[Description]: used by the \vpett{\nameref{mod:ComponentsPatientGatewayAppAppEndpoint}}.
      \item[Operations:] ~
    \begin{itemize}[noitemsep,nolistsep,leftmargin=-.25cm]
      \item \textsf{boolean pingConnection()}
        \begin{itemize}[noitemsep,nolistsep]
           \item Effect: In degraded mode, check for connection solved to restore buffer data. 
           \item Returns: Returns boolean-ping, indicating True for connection and lost message if failure.
           \item Sequence Diagrams: \cref{diag:Interaction:Ping}
        \end{itemize}
      \item \textsf{void processBufferedData(\ref{data:DatatypesSensorDataPackage} data)}
        \begin{itemize}[noitemsep,nolistsep]
           \item Effect:   Responsible for sending buffered data to PMs system after recovery. 
           \item Sequence Diagrams: \cref{diag:Interaction:Exitdegradedmode}
        \end{itemize}
    \end{itemize}
      \item[Diagrams:] None
    \end{description}

  %%%%%%%% SystemEventListener
  \subsection{SystemEventListener}\label{int:InterfacesSystemEventListener}
    \begin{description}[noitemsep,nolistsep]
      \item[Provided by:] \iconcomponent{}~\vpett{\nameref{mod:ComponentsPatientGatewayAppAppEndpoint}}, \iconcomponent{}~\vpett{\nameref{comp:ComponentsPatientGatewayApp}}
      \item[Required by:] \iconcomponent{}~\vpett{\nameref{comp:ComponentsMobileOS}} 
      \item[Operations:] ~
    \begin{itemize}[noitemsep,nolistsep,leftmargin=-.25cm]
      \item \textsf{void onBatteryLow()}
        \begin{itemize}[noitemsep,nolistsep]
           \item Effect: send an alert to the app that battery lv of the smartphone is under 15\% 
           \item Sequence Diagrams: \cref{diag:Interaction:Patientgatewayfails}
        \end{itemize}
      \item \textsf{void onBatteryVeryLow()}
        \begin{itemize}[noitemsep,nolistsep]
           \item Effect: send an alert to the app that the battery lvl of the smartphone is under 5\% 
           \item Sequence Diagrams: \cref{diag:Interaction:Patientgatewayfails}
        \end{itemize}
      \item \textsf{void onNetworkLost()}
        \begin{itemize}[noitemsep,nolistsep]
           \item Effect: send alert to the app that the network is lost 
           \item Sequence Diagrams: \cref{diag:Interaction:Patientgatewayfails}
        \end{itemize}
      \item \textsf{void onNetworkRestored()}
        \begin{itemize}[noitemsep,nolistsep]
           \item Effect: The mobile operating system lets the app now the network has restored. 
           \item Sequence Diagrams: None
        \end{itemize}
    \end{itemize}
      \item[Diagrams:] None
    \end{description}

  %%%%%%%% ThirdPartyAPIConnection
  \subsection{ThirdPartyAPIConnection}\label{int:InterfacesThirdPartyAPIConnection}
    \begin{description}[noitemsep,nolistsep]
      \item[Provided by:] \iconcomponent{}~\vpett{\nameref{mod:ComponentsThirdPartyServiceThirdPartyAPIsGoogleFitAPI}}, \iconcomponent{}~\vpett{\nameref{mod:ComponentsThirdPartyServiceThirdPartyAPIsStravaAPI}}, \iconcomponent{}~\vpett{\nameref{mod:ComponentsThirdPartyServiceThirdPartyAPIs}}
      \item[Required by:] \iconcomponent{}~\vpett{\nameref{mod:ComponentsThirdPartyServiceThirdPartySyncService}} 
           \item[Description]: an interface to connect with the api's
      \item[Operations:] ~
    \begin{itemize}[noitemsep,nolistsep,leftmargin=-.25cm]
      \item \textsf{\ref{data:DatatypesThirdPartyDataPackage} fetchData(\ref{data:DatatypesTimestamp} startTime, \ref{data:DatatypesTimestamp} endTime, string token)}
        \begin{itemize}[noitemsep,nolistsep]
           \item Effect: fetch the data 
           \item Returns: the data recieved from the api
           \item Sequence Diagrams: \cref{diag:Interaction:ThirdPartyNewData}
        \end{itemize}
    \end{itemize}
      \item[Diagrams:] None
    \end{description}

  %%%%%%%% ThirdPartyApiInterface
  \subsection{ThirdPartyApiInterface}\label{int:InterfacesThirdPartyApiInterface}
    \begin{description}[noitemsep,nolistsep]
      \item[Provided by:] \iconcomponent{}~\vpett{\nameref{comp:ComponentsThirdPartyService}}, \iconcomponent{}~\vpett{\nameref{mod:ComponentsThirdPartyServiceThirdPartySyncService}}
      \item[Required by:] \iconcomponent{}~\vpett{\nameref{comp:ComponentseHealthPlatformPatientDataStorage}}, \iconcomponent{}~\vpett{\nameref{comp:ComponentseHealthPlatform}} 
      \item[Operations:] ~
    \begin{itemize}[noitemsep,nolistsep,leftmargin=-.25cm]
      \item \textsf{\ref{data:DatatypesThirdPartyDataPackage} fetchData(\ref{data:DatatypesPatientId} patientId, \ref{data:DatatypesTimestamp} startTime, \ref{data:DatatypesTimestamp} endTime, String apiID)}
        \begin{itemize}[noitemsep,nolistsep]
           \item Effect: For fetching third party data from any service 
           \item Returns: Returns the third party data as a ThirdParty Data package
           \item Sequence Diagrams: \cref{diag:Interaction:Thirdpartypolling}
        \end{itemize}
    \end{itemize}
      \item[Diagrams:] None
    \end{description}

  %%%%%%%% ThirdPartyDataMgmt
  \subsection{ThirdPartyDataMgmt}\label{int:InterfacesThirdPartyDataMgmt}
    \begin{description}[noitemsep,nolistsep]
      \item[Provided by:] \iconcomponent{}~\vpett{\nameref{comp:ComponentseHealthPlatformPatientDataStorage}}
      \item[Required by:] \iconcomponent{}~\vpett{\nameref{comp:ComponentseHealthPlatformClinicalJobCreator}} 
           \item[Description]: Interface to the database storing data from third party devices
      \item[Operations:] ~
    \begin{itemize}[noitemsep,nolistsep,leftmargin=-.25cm]
      \item \textsf{Map\textless{}\ref{data:DatatypesTimestamp}, \ref{data:DatatypesThirdPartyDataPackage}\textgreater{} getTPPatientData(\ref{data:DatatypesPatientId} patient, \ref{data:DatatypesTimestamp} startDate, \ref{data:DatatypesTimestamp} endDate, String Api)}
        \begin{itemize}[noitemsep,nolistsep]
           \item Effect: For fetching the data of the third party services 
           \item Returns: returns the third party patient data
           \item Sequence Diagrams: \cref{diag:Interaction:Thirdpartypolling}
        \end{itemize}
    \end{itemize}
      \item[Diagrams:] None
    \end{description}

  %%%%%%%% TokenMgmt
  \subsection{TokenMgmt}\label{int:InterfacesTokenMgmt}
    \begin{description}[noitemsep,nolistsep]
      \item[Provided by:] \iconcomponent{}~\vpett{\nameref{comp:ComponentsThirdPartyService}}, \iconcomponent{}~\vpett{\nameref{mod:ComponentsThirdPartyServiceTokenStorage}}
      \item[Required by:] \iconcomponent{}~\vpett{\nameref{comp:ComponentsPatientGatewayEndpoint}}, \iconcomponent{}~\vpett{\nameref{mod:ComponentsThirdPartyServiceThirdPartySyncService}} 
      \item[Operations:] ~
    \begin{itemize}[noitemsep,nolistsep,leftmargin=-.25cm]
      \item \textsf{string getToken(\ref{data:DatatypesPatientId} patientId, string provider)}
        \begin{itemize}[noitemsep,nolistsep]
           \item Effect: looks for a specific token in the database and returns it, so it can be used for the third party api 
           \item Returns: the token as a string
           \item Sequence Diagrams: \cref{diag:Interaction:ThirdPartyNewData}
        \end{itemize}
      \item \textsf{void saveNewTokens(List\textless{}\ref{data:DatatypesTokenPackage}\textgreater{} newTokens)}
        \begin{itemize}[noitemsep,nolistsep]
           \item Effect: saves the tokens in the database based on the patientId, provider and the token itself 
           \item Sequence Diagrams: \cref{diag:Interaction:Storetokens}
        \end{itemize}
    \end{itemize}
      \item[Diagrams:] None
    \end{description}

% END INTERFACES

% NODES
\section{Nodes}\label{sec:nodes}
\subsection{ControlNode}\label{node:DeploymentControlNode}
	\begin{description}[noitemsep,nolistsep]
		\item[Responsibility:]~The node for controlling the sensordata fetches using the endpoints, \vpett{\nameref{comp:ComponentseHealthPlatformHIS}} and chronjob.
		\item[Visible on diagrams:]~\cref{diag:Deployment:DevelopmentTestDeployment20patients3clinmodels,diag:Deployment:PilotDeployment200patients5clinmodels,diag:Deployment:Primarydeploymentdiagram}		
	\end{description}

\subsection{DBNode (Pilot Deployment)}\label{node:DeploymentDBNodePilotDeployment}
	\begin{description}[noitemsep,nolistsep]
		\item[Responsibility:]~node for db storage and fetching third party data
		\item[Visible on diagrams:]~\cref{diag:Deployment:PilotDeployment200patients5clinmodels}		
	\end{description}

\subsection{eHealthPlatform}\label{node:DeploymenteHealthPlatform}
	\begin{description}[noitemsep,nolistsep]
		\item[Responsibility:]~Container node representing the deployment context boundary for the \vpett{\nameref{comp:ComponentseHealthPlatform}}.
		\item[Visible on diagrams:]~\cref{diag:Deployment:Contextdiagramforthedeploymentview}		
	\end{description}

\subsection{MLaaS Conn (Pilot Deployment)}\label{node:DeploymentMLaaSConnPilotDeployment}
	\begin{description}[noitemsep,nolistsep]
		\item[Responsibility:]~This node is responsible for the connection with the \vpett{\nameref{comp:ComponentsMLaaS}} Service.
		\item[Visible on diagrams:]~\cref{diag:Deployment:PilotDeployment200patients5clinmodels}		
	\end{description}

\subsection{MLaaS Service (Dev Deployment)}\label{node:DeploymentMLaaSServiceDevDeployment}
	\begin{description}[noitemsep,nolistsep]
		\item[Responsibility:]~This represents the \vpett{\nameref{comp:ComponentsMLaaS}} cloud service.
		\item[Visible on diagrams:]~\cref{diag:Deployment:DevelopmentTestDeployment20patients3clinmodels}		
	\end{description}

\subsection{MLaaS Service (Pilot Deployment)}\label{node:DeploymentMLaaSServicePilotDeployment}
	\begin{description}[noitemsep,nolistsep]
		\item[Responsibility:]~This represents the cloud \vpett{\nameref{comp:ComponentsMLaaS}} service.
		\item[Visible on diagrams:]~\cref{diag:Deployment:PilotDeployment200patients5clinmodels}		
	\end{description}

\subsection{MLaaSCloud}\label{node:DeploymentMLaaSCloud}
	\begin{description}[noitemsep,nolistsep]
		\item[Responsibility:]~Cloud infrastructure of the \vpett{\nameref{comp:ComponentsMLaaS}} Service provider.
		\item[Visible on diagrams:]~\cref{diag:Deployment:Contextdiagramforthedeploymentview}		
	\end{description}

\subsection{MLaaSConnectorNode}\label{node:DeploymenteHealthPlatformMLaaSConnectorNode}
	\begin{description}[noitemsep,nolistsep]
		\item[Responsibility:]~Node responsible for handling connections with external \vpett{\nameref{comp:ComponentsMLaaS}} providers.
		\item[Visible on diagrams:]~\cref{diag:Deployment:Contextdiagramforthedeploymentview,diag:Deployment:Primarydeploymentdiagram}		
	\end{description}

\subsection{modelDB (Pilot Deployment)}\label{node:DeploymentmodelDBPilotDeployment}
	\begin{description}[noitemsep,nolistsep]
		\item[Responsibility:]~extra node for storing the models 
		\item[Visible on diagrams:]~\cref{diag:Deployment:PilotDeployment200patients5clinmodels}		
	\end{description}

\subsection{ModelStorage}\label{node:DeploymentModelStorage}
	\begin{description}[noitemsep,nolistsep]
		\item[Responsibility:]~For storing and providing the ml models
		\item[Visible on diagrams:]~\cref{diag:Deployment:Primarydeploymentdiagram}		
	\end{description}

\subsection{Patient Devices}\label{node:DeploymentPatientDevices}
	\begin{description}[noitemsep,nolistsep]
		\item[Responsibility:]~the devices (smartphones) of the patients

		\item[Visible on diagrams:]~\cref{diag:Deployment:Primarydeploymentdiagram}		
	\end{description}

\subsection{PatientDB}\label{node:DeploymentPatientDB}
	\begin{description}[noitemsep,nolistsep]
		\item[Responsibility:]~Node for storing the third party patient data and fetching the third party data from the services
		\item[Visible on diagrams:]~\cref{diag:Deployment:DevelopmentTestDeployment20patients3clinmodels,diag:Deployment:Primarydeploymentdiagram}		
	\end{description}

\subsection{RiskEstimationManagementNode (Pilot Deployment)}\label{node:DeploymentRiskEstimationManagementNodePilotDeployment}
	\begin{description}[noitemsep,nolistsep]
		\item[Responsibility:]~The RiskEstimationManagement node hosts the creation, scheduling and combining of the clinical jobs.
		\item[Visible on diagrams:]~\cref{diag:Deployment:PilotDeployment200patients5clinmodels}		
	\end{description}

\subsection{RiskEstimationManagmentNode (Dev Deployment)}\label{node:DeploymentRiskEstimationManagmentNodeDevDeployment}
	\begin{description}[noitemsep,nolistsep]
		\item[Responsibility:]~The development deployment is a single RiskEstimationManagement node that contains all other functionality, job creation, scheduling, and combination.
		\item[Visible on diagrams:]~\cref{diag:Deployment:DevelopmentTestDeployment20patients3clinmodels}		
	\end{description}

\subsection{RiskEstimationMgmtNode}\label{node:DeploymenteHealthPlatformRiskEstimationMgmtNode}
	\begin{description}[noitemsep,nolistsep]
		\item[Responsibility:]~The RiskEstimationMgmtNode keeps track of the risk estimation jobs, and schedules and combines the results.
		\item[Visible on diagrams:]~\cref{diag:Deployment:Primarydeploymentdiagram}		
	\end{description}

\subsection{RiskEstimationProcessorNode}\label{node:DeploymenteHealthPlatformRiskEstimationProcessorNode}
	\begin{description}[noitemsep,nolistsep]
		\item[Responsibility:]~Multiple \vpett{\nameref{comp:ComponentseHealthPlatformRiskEstimationProcessor}} nodes enable the processing of multiple sensorData inputs in parallel.
		\item[Visible on diagrams:]~\cref{diag:Deployment:Primarydeploymentdiagram}		
	\end{description}

\subsection{RiskEstimationProcessorNode (Dev Deployment)}\label{node:DeploymentRiskEstimationProcessorNodeDevDeployment}
	\begin{description}[noitemsep,nolistsep]
		\item[Responsibility:]~The development deployment only uses two RiskEstimationProcessorNodes to run \ref{data:DatatypesClinicalModelJob}s in parallel.
		\item[Visible on diagrams:]~\cref{diag:Deployment:DevelopmentTestDeployment20patients3clinmodels}		
	\end{description}

\subsection{RiskEstimationProcessorNode (Pilot Deployment)}\label{node:DeploymentRiskEstimationProcessorNodePilotDeployment}
	\begin{description}[noitemsep,nolistsep]
		\item[Responsibility:]~There are five instances of the RiskEstimationProcessorNode to enable the processing of multiple \ref{data:DatatypesClinicalModelJob}s in parallel.
		\item[Visible on diagrams:]~\cref{diag:Deployment:PilotDeployment200patients5clinmodels}		
	\end{description}

\subsection{SMS Provider}\label{node:DeploymenteHealthPlatformSMSProvider}
	\begin{description}[noitemsep,nolistsep]
		\item[Responsibility:]~Node responsible for handling connections with external \vpett{\nameref{comp:ComponentsMLaaS}} providers.
		\item[Visible on diagrams:]~\cref{diag:Deployment:Primarydeploymentdiagram}		
	\end{description}


% END NODES

% EXCEPTIONS
\section{Exceptions}\label{sec:exceptions}
\begin{itemize}[nolistsep,noitemsep]
\item \vpeexception{ex:ExceptionsAuthenticationException}{AuthenticationException} (\iconinher{}~\ref{ex:ExceptionsException}): 
\begin{itemize}[noitemsep,nolistsep]

\item[] Signals a problem with the credentials used for authentication. These may be invalid or expired.
\end{itemize}

\item \vpeexception{ex:ExceptionsAuthorizationException}{AuthorizationException} (\iconinher{}~\ref{ex:ExceptionsException}): 
\begin{itemize}[noitemsep,nolistsep]

\item[] Signals that the principal has insufficient access rights to perform the operation.
\end{itemize}

\item \vpeexception{ex:ExceptionsException}{Exception}: 
\begin{itemize}[noitemsep,nolistsep]
\item[] Subtypes: \iconinher{}~\ref{ex:ExceptionsAuthenticationException}, \iconinher{}~\ref{ex:ExceptionsAuthorizationException}, \iconinher{}~\ref{ex:ExceptionsInvalidAPIKeyException}, \iconinher{}~\ref{ex:ExceptionsInvalidMLModelException}, \iconinher{}~\ref{ex:ExceptionsIOException}, \iconinher{}~\ref{ex:ExceptionsNoSuchMLModelException}, \iconinher{}~\ref{ex:ExceptionsNoSuchPatientException}, \iconinher{}~\ref{ex:ExceptionsNoSuchPatientRecordException}
\item[] Base exception class.
\end{itemize}

\item \vpeexception{ex:ExceptionsInvalidAPIKeyException}{InvalidAPIKeyException} (\iconinher{}~\ref{ex:ExceptionsException}): 
\begin{itemize}[noitemsep,nolistsep]

\item[] Thrown when no correctly formatted API key is configured for constructing authentication Credential|s.
\end{itemize}

\item \vpeexception{ex:ExceptionsInvalidMLModelException}{InvalidMLModelException} (\iconinher{}~\ref{ex:ExceptionsException}): 
\begin{itemize}[noitemsep,nolistsep]

\item[] Thrown when a provided MLModel is invalid.
\end{itemize}

\item \vpeexception{ex:ExceptionsIOException}{IOException} (\iconinher{}~\ref{ex:ExceptionsException}): 
\begin{itemize}[noitemsep,nolistsep]

\item[] Thrown when an IO operation fails.
\end{itemize}

\item \vpeexception{ex:ExceptionsNoSuchMLModelException}{NoSuchMLModelException} (\iconinher{}~\ref{ex:ExceptionsException}): 
\begin{itemize}[noitemsep,nolistsep]

\item[] Thrown when the requested MLModel cannot be found.
\end{itemize}

\item \vpeexception{ex:ExceptionsNoSuchPatientException}{NoSuchPatientException} (\iconinher{}~\ref{ex:ExceptionsException}): 
\begin{itemize}[noitemsep,nolistsep]

\item[] Thrown if no patient with the specified identifier exists.
\end{itemize}

\item \vpeexception{ex:ExceptionsNoSuchPatientRecordException}{NoSuchPatientRecordException} (\iconinher{}~\ref{ex:ExceptionsException}): 
\begin{itemize}[noitemsep,nolistsep]

\item[] Thrown if no patient record for the patient with the given identifier exists in the HIS.
\end{itemize}

\end{itemize}
% END EXCEPTIONS

% DATA TYPES
\section{Data types}\label{sec:datatypes}
\begin{itemize}[nolistsep,noitemsep]
\item \vpedatatype{data:Datatypesorghl7fhirr5modelAddress}{Address}: 
\begin{itemize}[noitemsep,nolistsep]
\item[] Attributes: \ref{data:Datatypesorghl7fhirr5modelAddressAddressUse} use, \ref{data:Datatypesorghl7fhirr5modelAddressAddressType} type, String text, List\textless{}String\textgreater{} line, String city, String district, String state, String postalCode, String country, \ref{data:Datatypesorghl7fhirr5modelPeriod} period, long serialVersionUID
\item[] Base StructureDefinition for Address Type: An address expressed using postal conventions (as opposed to GPS or other location definition formats).  This data type may be used to convey addresses for use in delivering mail as well as for visiting locations which might not be valid for mail delivery.  There are a variety of postal address formats defined around the world.
The line attribute contains contains the house number, apartment number, street name, street direction, P.O.
The text attribute specifies the entire address as it should be displayed e.g. on a postal label. This may be provided instead of or as well as the specific parts.
\end{itemize}

\item \vpedatatype{data:Datatypesorghl7fhirr5modelAddressAddressType}{AddressType}: 
\begin{itemize}[noitemsep,nolistsep]

\item[] Enum: \textsc{postal}, \textsc{physical}, \textsc{both}, \textsc{null}.\\Distinguishes between physical addresses (those you can visit) and mailing addresses (e.g. PO Boxes and care-of addresses). Most addresses are both.
\end{itemize}

\item \vpedatatype{data:Datatypesorghl7fhirr5modelAddressAddressUse}{AddressUse}: 
\begin{itemize}[noitemsep,nolistsep]

\item[] Enum: \textsc{home}, \textsc{work}, \textsc{temp}, \textsc{old}, \textsc{billing}, \textsc{null}.\\The purpose of the address.
\end{itemize}

\item \vpedatatype{data:Datatypesorghl7fhirr5modelAnnotation}{Annotation}: 
\begin{itemize}[noitemsep,nolistsep]
\item[] Attributes: \ref{data:DatatypesDate} time, String text, long serialVersionUID, String author
\item[] Base StructureDefinition for Annotation Type: A  text note which also  contains information about who made the statement and when.
\end{itemize}

\item \vpedatatype{data:DatatypesBigDecimal}{BigDecimal}: 
\begin{itemize}[noitemsep,nolistsep]

\item[] Immutable, arbitrary-precision signed decimal numbers. A BigDecimal consists of an arbitrary precision integer unscaled value and a 32-bit integer scale
\end{itemize}

\item \vpedatatype{data:DatatypesClinicalModel}{ClinicalModel}: 
\begin{itemize}[noitemsep,nolistsep]
\item[] Attributes: String id, String name, String description, List\textless{}String\textgreater{} requiredParameters, List\textless{}String\textgreater{} optionalParameters, List\textless{}String\textgreater{} applicableMLModelIds, boolean requiresHistoricalSensorData
\item[] A clinical model. This is model has an id, name, description.
It contains a list of  parameters for using the model and a set of MLModelIds to specify which types of MLModels can be used in the context of the ClinicalModel.

\end{itemize}

\item \vpedatatype{data:DatatypesClinicalModelConfiguration}{ClinicalModelConfiguration}: 
\begin{itemize}[noitemsep,nolistsep]
\item[] Attributes: Map\textless{}String, String\textgreater{} modelParameters
\item[] The configuration for a clinical model. Represented as a set key/value pair for the parameters of the clinical model.
\end{itemize}

\item \vpedatatype{data:DatatypesClinicalModelJob}{ClinicalModelJob}: 
\begin{itemize}[noitemsep,nolistsep]
\item[] Attributes: \ref{data:DatatypesClinicalModel} clinicalModel, \ref{data:DatatypesClinicalModelJobId} jobId, \ref{data:DatatypesSensorDataPackage} sensorData, \ref{data:DatatypesPatientId} patientId, \ref{data:DatatypesClinicalModelConfiguration} clinicalModelConfiguration, \ref{data:DatatypesPatientRecord} patientRecord, Map\textless{}\ref{data:DatatypesTimestamp}, \ref{data:DatatypesSensorDataPackage}\textgreater{} historicalSensorData, \ref{data:DatatypesThirdPartyDataPackage} thirdPartyData, Map\textless{}\ref{data:DatatypesTimestamp}, \ref{data:DatatypesThirdPartyDataPackage}\textgreater{} historicalThirdPartyData
\item[] A single job for the computation of a clinical model. This contains a ClinicalModelJobId, ClinicalModelId and PatientId respectively identifying the job itself, the corresponding clinical model and the patient. If the corresponding risk estimation the job belongs to is triggered by the arrival of a sensor data package, this sensor data is also contained.
\end{itemize}

\item \vpedatatype{data:DatatypesClinicalModelJobId}{ClinicalModelJobId}: 
\begin{itemize}[noitemsep,nolistsep]

\item[]  A piece of data uniquely identifying a certain clinical model job in the system. This architecture does not specify the exact format of this identifier,
but possibilities are a long integer, a string, a URI etc.
\end{itemize}

\item \vpedatatype{data:DatatypesClinicalModelJobResult}{ClinicalModelJobResult}: 
\begin{itemize}[noitemsep,nolistsep]

\item[] The result of the computation of a clinical model containing the estimated
risk level for the patient.
The resulting risk score for ClinicalModels are on the same unit scale so the the results of multiple ClinicalModels can be combined.
\end{itemize}

\item \vpedatatype{data:Datatypesorghl7fhirr5modelCodeableConcept}{CodeableConcept}: 
\begin{itemize}[noitemsep,nolistsep]
\item[] Attributes: List\textless{}\ref{data:Datatypesorghl7fhirr5modelCoding}\textgreater{} coding, String text, long serialVersionUID
\item[] Base StructureDefinition for CodeableConcept Type: A concept that may be defined by a formal reference to a terminology or ontology or may be provided by text.
\end{itemize}

\item \vpedatatype{data:Datatypesorghl7fhirr5modelCodeType}{CodeType}: 
\begin{itemize}[noitemsep,nolistsep]
\item[] Attributes: long serialVersionUID, String system, String value
\item[] Primitive type "code" in FHIR, when not bound to an enumerated list of codes
\end{itemize}

\item \vpedatatype{data:Datatypesorghl7fhirr5modelCoding}{Coding}: 
\begin{itemize}[noitemsep,nolistsep]
\item[] Attributes: String system, String version, \ref{data:Datatypesorghl7fhirr5modelCodeType} code, String display, boolean userSelected, long serialVersionUID
\item[] Base StructureDefinition for Coding Type: A reference to a code defined by a terminology system.
\end{itemize}

\item \vpedatatype{data:Datatypesorghl7fhirr5modelPatientContactComponent}{ContactComponent}: 
\begin{itemize}[noitemsep,nolistsep]
\item[] Attributes: List\textless{}\ref{data:Datatypesorghl7fhirr5modelCodeableConcept}\textgreater{} relationship, String name, List\textless{}\ref{data:Datatypesorghl7fhirr5modelContactPoint}\textgreater{} telecom, \ref{data:Datatypesorghl7fhirr5modelAddress} address, String administrativeGender, String organization, \ref{data:Datatypesorghl7fhirr5modelPeriod} period, long serialVersionUID
\item[] A contact party (e.g. guardian, partner, friend) for the patient.
The relationship attribute species the nature of the relationship between the patient and the contact person.
It is a CodeableConcept (which can be formal reference to terminology or an ontology).
For example, in the system \url{http://hl7.org/fhir/ValueSet/relatedperson-relationshiptype}, it can be coded as CHILD for a child, or MGRMTH for maternal grandmother, etc.
\end{itemize}

\item \vpedatatype{data:Datatypesorghl7fhirr5modelContactPoint}{ContactPoint}: 
\begin{itemize}[noitemsep,nolistsep]
\item[] Attributes: \ref{data:Datatypesorghl7fhirr5modelContactPointContactPointSystem} system, String value, \ref{data:Datatypesorghl7fhirr5modelContactPointContactPointUse} use, int rank, long serialVersionUID, \ref{data:Datatypesorghl7fhirr5modelPeriod} period
\item[] Base StructureDefinition for ContactPoint Type: Details for all kinds of technology mediated contact points for a person or organization, including telephone, email, etc.
The value attribute contains the actual contact point details, in a form that is meaningful to the designated communication system (i.e. phone number or email address).
Rank specifies a preferred order in which to use a set of contacts. ContactPoints with lower rank values are more preferred than those with higher rank values.
Period specifies the tiime period when the contact point was/is in use.

\end{itemize}

\item \vpedatatype{data:Datatypesorghl7fhirr5modelContactPointContactPointSystem}{ContactPointSystem}: 
\begin{itemize}[noitemsep,nolistsep]

\item[] Enum: \textsc{phone}, \textsc{fax}, \textsc{email}, \textsc{pager}, \textsc{url}, \textsc{sms}, \textsc{other}, \textsc{null}.\\Enum denotying the type of communication channel for the ContactPoint.
\end{itemize}

\item \vpedatatype{data:Datatypesorghl7fhirr5modelContactPointContactPointUse}{ContactPointUse}: 
\begin{itemize}[noitemsep,nolistsep]

\item[] Enum: \textsc{home}, \textsc{work}, \textsc{temp}, \textsc{old}, \textsc{mobile}, \textsc{null}.\\Enum denoting the usage type of the ContactPoint.
\end{itemize}

\item \vpedatatype{data:DatatypesCredential}{Credential}: 
\begin{itemize}[noitemsep,nolistsep]

\item[] Authentication object for making authenticated requests to the MLaaS service provider.
\end{itemize}

\item \vpedatatype{data:DatatypesDate}{Date}: 
\begin{itemize}[noitemsep,nolistsep]

\item[] The class Date represents a specific instant in time, with millisecond precision. 
\end{itemize}

\item \vpedatatype{data:Datatypesorghl7fhirr5modelIdentifier}{Identifier}: 
\begin{itemize}[noitemsep,nolistsep]
\item[] Attributes: \ref{data:Datatypesorghl7fhirr5modelIdentifierIdentifierUse} use, String type, String system, String value, \ref{data:Datatypesorghl7fhirr5modelPeriod} period, String assigner, long serialVersionUID
\item[] Base StructureDefinition for Identifier Type: An identifier - identifies some entity uniquely and unambiguously. Typically this is used for business identifiers.
The type is a a coded type for the identifier that can be used to determine which identifier to use for a specific purpose.
The system attribute establishes the namespace for the value - that is, a URL that describes a set values that are unique.
The value attribute specifies the portion of the identifier typically relevant to the user and which is unique within the context of the system.
The period specifies the time period during which identifier is/was valid for use.
The assigner specifies the organization that issued/manages the identifier.
\end{itemize}

\item \vpedatatype{data:Datatypesorghl7fhirr5modelIdentifierIdentifierUse}{IdentifierUse}: 
\begin{itemize}[noitemsep,nolistsep]

\item[] Enum: \textsc{usual}, \textsc{official}, \textsc{temp}, \textsc{secondary}, \textsc{old}, \textsc{null}.\\Enum denoting the usage type of a particular Identifier.
\end{itemize}

\item \vpedatatype{data:DatatypesMLModel}{MLModel}: 
\begin{itemize}[noitemsep,nolistsep]
\item[] Attributes: String modelId, \ref{data:DatatypesMLModelDescription} modelDescription
\item[] Machine learning model for making predictions based on sensordata.
The actual internal encoding of the model itself is left open.
\end{itemize}

\item \vpedatatype{data:DatatypesMLModelDescription}{MLModelDescription}: 
\begin{itemize}[noitemsep,nolistsep]
\item[] Attributes: String uuid, String name, String description, String providerName
\item[] Object providing the MLModel meta data. It does not contain the model itself.
\end{itemize}

\item \vpedatatype{data:DatatypesMLModelUpdate}{MLModelUpdate}: 
\begin{itemize}[noitemsep,nolistsep]
\item[] Attributes: \ref{data:DatatypesMLModelDescription} modelDescription
\item[] Represents a model update that can be applied on an MLModel.
These model updates enable exchanging changes to the MLModel between different parties without sending the actual MLModel to the other parties.
\end{itemize}

\item \vpedatatype{data:DatatypesObject}{Object}: 
\begin{itemize}[noitemsep,nolistsep]

\item[] Root class for objects that can be stored in the Key-Value storage service.
Data types that need to be saved need to inherit from this class.
\end{itemize}

\item \vpedatatype{data:Datatypesorghl7fhirr5modelObservation}{Observation}: 
\begin{itemize}[noitemsep,nolistsep]
\item[] Attributes: List\textless{}\ref{data:Datatypesorghl7fhirr5modelIdentifier}\textgreater{} identifier, List\textless{}\ref{data:Datatypesorghl7fhirr5modelIdentifier}\textgreater{} partOf, \ref{data:Datatypesorghl7fhirr5modelObservationStatus} status, List\textless{}\ref{data:Datatypesorghl7fhirr5modelCodeableConcept}\textgreater{} category, \ref{data:Datatypesorghl7fhirr5modelCodeableConcept} code, \ref{data:Datatypesorghl7fhirr5modelPatient} subject, \ref{data:Datatypesorghl7fhirr5modelPeriod} effective, \ref{data:DatatypesDate} issued, List\textless{}\ref{data:Datatypesorghl7fhirr5modelIdentifier}\textgreater{} performer, String value, \ref{data:Datatypesorghl7fhirr5modelCodeableConcept} dataAbsentReason, \ref{data:Datatypesorghl7fhirr5modelCodeableConcept} interpretation, List\textless{}\ref{data:Datatypesorghl7fhirr5modelAnnotation}\textgreater{} note, \ref{data:Datatypesorghl7fhirr5modelCodeableConcept} bodySite, \ref{data:Datatypesorghl7fhirr5modelCodeableConcept} method, String device, List\textless{}\ref{data:Datatypesorghl7fhirr5modelObservation}\textgreater{} hasMember, List\textless{}\ref{data:Datatypesorghl7fhirr5modelObservation}\textgreater{} derivedFrom, long serialVersionUID, \ref{data:Datatypesorghl7fhirr5modelObservationObservationReferenceRangeComponent} referenceRange, \ref{data:Datatypesorghl7fhirr5modelObservationObservationComponentComponent} component
\item[] Measurements and simple assertions made about a patient, device or other subject.
The partOf attribute specifies the larger event of which this particular Observation is a component or step. For example, an observation as part of a procedure.

The category specifies a code that classifies the general type of observation being made.
For example, in the system \url{http://terminology.hl7.org/CodeSystem/observation-category}, vital-signs is used for clinical observations such as boold pressure, heart rate, etc.

The code attribute describes what was observed. Sometimes this is called the observation `name'.
For example, in the system \url{http://loinc.org}, code `8867-4' is used for expressing the heart rate, `85354-9' for a blood pressure panel (which must have components `8480-6' and `8462-4' for, respectively, systolic and diastolic pressure.

The performer is who was responsible for asserting the observed value as `true'. This can refer to a Patient, Practitioner, organization, etc.
The value is the information determined as a result of making the observation, if the information has a simple value.

The dataAbsentReason provides a reason why the expected value in the element Observation.value[x] is missing.
For example, in the system \url{http://terminology.hl7.org/CodeSystem/data-absent-reason}, `not-applicable'.

The interpretation is a categorical assessment of an observation value.  For example, high, low, normal.
For example, in the system \url{http://terminology.hl7.org/CodeSystem/v3-ObservationInterpretation}, `H' for high, or `ND' for not detected.

The bodySite indicates the site on the subject's body where the observation was made (i.e. the target site).
For example, in the system \url{http://snomed.info/sct}, code `344001' indicates the ankle.

The method indicates the mechanism used to perform the observation.
For example, in the system \url{http://snomed.info/sct}, code `46973005' indicates blood pressure taking procedure.

The device attribute specifies an identifier of the device used to generate the observation data.
\end{itemize}

\item \vpedatatype{data:Datatypesorghl7fhirr5modelObservationObservationComponentComponent}{ObservationComponentComponent}: 
\begin{itemize}[noitemsep,nolistsep]
\item[] Attributes: String code, \ref{data:Datatypesorghl7fhirr5modelRange} value, String dataAbsentReason, List\textless{}String\textgreater{} interpretation, long serialVersionUID, \ref{data:Datatypesorghl7fhirr5modelObservationObservationReferenceRangeComponent} referenceRange
\item[] Some observations have multiple component observations.  These component observations are expressed as separate code value pairs that share the same attributes.  Examples include systolic and diastolic component observations for blood pressure measurement and multiple component observations for genetics observations.
The code describes what was observed. Sometimes this is called the observation `code'.
For example, in the system \url{http://loinc.org}, `8480-6' is used for the systolic component of the blood pressure.,

The dataAbsentReason provides a reason why the expected value in the element Observation.component.value[x] is missing.
For example, in the system \url{http://terminology.hl7.org/CodeSystem/data-absent-reason}, `not-applicable'.

The interpretation provides a categorical assessment of an observation value. For example, high, low, normal.
For example, in the system \url{http://terminology.hl7.org/CodeSystem/v3-ObservationInterpretation}, `H' for high, or `ND' for not detected.
\end{itemize}

\item \vpedatatype{data:Datatypesorghl7fhirr5modelObservationObservationReferenceRangeComponent}{ObservationReferenceRangeComponent}: 
\begin{itemize}[noitemsep,nolistsep]
\item[] Attributes: \ref{data:Datatypesorghl7fhirr5modelQuantity} low, \ref{data:Datatypesorghl7fhirr5modelQuantity} high, \ref{data:Datatypesorghl7fhirr5modelCodeableConcept} type, List\textless{}\ref{data:Datatypesorghl7fhirr5modelCodeableConcept}\textgreater{} appliesTo, \ref{data:Datatypesorghl7fhirr5modelRange} age, String text, long serialVersionUID
\item[] Guidance on how to interpret the value by comparison to a normal or recommended range.  Multiple reference ranges are interpreted as an "OR".   In other words, to represent two distinct target populations, two `referenceRange` elements would be used.

The appliesTo specifies codes to indicate the target population this reference range applies to.
For example, in the system \url{http://snomed.info/sct}, code `248152002' indicates female.

The type specifies codes to indicate the what part of the targeted reference population it applies to.
For example, in the system \url{http://terminology.hl7.org/CodeSystem/referencerange-meaning}, `normal' for the 95\% normal range, or `recommended' for the recommended range by a relevant professional body.

The text attribute contains a text based reference range in an observation which may be used when a quantitative range is not appropriate for an observation.
\end{itemize}

\item \vpedatatype{data:Datatypesorghl7fhirr5modelObservationStatus}{ObservationStatus}: 
\begin{itemize}[noitemsep,nolistsep]

\item[] Enum: \textsc{registered}, \textsc{preliminary}, \textsc{final}, \textsc{amended}, \textsc{corrected}, \textsc{cancelled}, \textsc{enteredinerror}, \textsc{unknown}, \textsc{null}.\\The status of the observation result value or the RiskAssessment.
\end{itemize}

\item \vpedatatype{data:Datatypesorghl7fhirr5modelPatient}{Patient}: 
\begin{itemize}[noitemsep,nolistsep]
\item[] Attributes: List\textless{}\ref{data:Datatypesorghl7fhirr5modelIdentifier}\textgreater{} identifier, boolean active, List\textless{}String\textgreater{} name, String administrativeGender, \ref{data:DatatypesDate} birthDate, \ref{data:DatatypesDate} deceased, List\textless{}\ref{data:Datatypesorghl7fhirr5modelAddress}\textgreater{} address, boolean multipleBirth, List\textless{}\ref{data:Datatypesorghl7fhirr5modelIdentifier}\textgreater{} generalPractitioner, \ref{data:Datatypesorghl7fhirr5modelIdentifier} managingOrganization, long serialVersionUID, \ref{data:Datatypesorghl7fhirr5modelPatientContactComponent} contact
\item[] Demographics and other administrative information about an individual or animal receiving care or other health-related services.
\end{itemize}

\item \vpedatatype{data:DatatypesPatientId}{PatientId}: 
\begin{itemize}[noitemsep,nolistsep]

\item[] A piece of data uniquely identifying a patient in the system. This architecture does not
specify the exact format of this identifier, but possibilities are a long integer, a string, a URI, etc.
\end{itemize}

\item \vpedatatype{data:DatatypesPatientRecord}{PatientRecord}: 
\begin{itemize}[noitemsep,nolistsep]

\item[] The structured contents of the EHR record of a certain patient, including medication history, recent treatments, allergies, etc.
\end{itemize}

\item \vpedatatype{data:DatatypesPatientStatus}{PatientStatus}: 
\begin{itemize}[noitemsep,nolistsep]

\item[] The status of a patient, i.e., their risk level.
\end{itemize}

\item \vpedatatype{data:Datatypesorghl7fhirr5modelPeriod}{Period}: 
\begin{itemize}[noitemsep,nolistsep]
\item[] Attributes: \ref{data:DatatypesDate} start, \ref{data:DatatypesDate} end, long serialVersionUID
\item[] Base StructureDefinition for Period Type: A time period defined by a start and end date and optionally time.
\end{itemize}

\item \vpedatatype{data:DatatypesPrediction}{Prediction}: 
\begin{itemize}[noitemsep,nolistsep]

\item[] A prediction from the MLModel.
\end{itemize}

\item \vpedatatype{data:Datatypesorghl7fhirr5modelQuantity}{Quantity}: 
\begin{itemize}[noitemsep,nolistsep]
\item[] Attributes: \ref{data:DatatypesBigDecimal} value, String unit, String system, \ref{data:Datatypesorghl7fhirr5modelCodeType} code, long serialVersionUID
\item[] Base StructureDefinition for Quantity Type: A measured amount (or an amount that can potentially be measured). Note that measured amounts include amounts that are not precisely quantified, including amounts involving arbitrary units and floating currencies.
Unit specifies a human-readable form of the unit.
System contains the identification of the system that provides the coded form of the unit.
The code contains a computer processable form of the unit in some unit representation system.
For example, blood pressure can be expressed in the system \url{http://unitsofmeasure.org} and the unit mm[Hg].
\end{itemize}

\item \vpedatatype{data:DatatypesQuery}{Query}\textbf{\textless{}T\textgreater{}}: 
\begin{itemize}[noitemsep,nolistsep]
\item[] Attributes: List\textless{}T\textgreater{} values, List\textless{}String\textgreater{} qualifier
\item[] The Query object is used to express search restrictions in a string-based format for interacting with the REST-based health API.
For example, before a specified date is encoded as `ltyyyy-MM-dd'.
The list of qualifiers specify how the corresponding value should be used (e.g., lt for less than a date, matches if it should be present in a string, exact for an exact string match).
\end{itemize}

\item \vpedatatype{data:Datatypesorghl7fhirr5modelRange}{Range}: 
\begin{itemize}[noitemsep,nolistsep]
\item[] Attributes: long serialVersionUID, \ref{data:Datatypesorghl7fhirr5modelQuantity} low, \ref{data:Datatypesorghl7fhirr5modelQuantity} high
\item[] Base StructureDefinition for Range Type: A set of ordered Quantities defined by a low and high limit.
\end{itemize}

\item \vpedatatype{data:Datatypesorghl7fhirr5modelRiskAssessment}{RiskAssessment}: 
\begin{itemize}[noitemsep,nolistsep]
\item[] Attributes: List\textless{}\ref{data:Datatypesorghl7fhirr5modelIdentifier}\textgreater{} identifier, \ref{data:Datatypesorghl7fhirr5modelObservationStatus} status, \ref{data:Datatypesorghl7fhirr5modelCodeableConcept} method, \ref{data:Datatypesorghl7fhirr5modelCodeableConcept} code, \ref{data:Datatypesorghl7fhirr5modelPatient} subject, \ref{data:DatatypesDate} occurrence, String condition, \ref{data:Datatypesorghl7fhirr5modelIdentifier} performer, List\textless{}\ref{data:Datatypesorghl7fhirr5modelCodeableConcept}\textgreater{} reason, List\textless{}\ref{data:Datatypesorghl7fhirr5modelObservation}\textgreater{} basis, String mitigation, List\textless{}\ref{data:Datatypesorghl7fhirr5modelAnnotation}\textgreater{} note, long serialVersionUID, \ref{data:Datatypesorghl7fhirr5modelRiskAssessmentRiskAssessmentPredictionComponent} prediction
\item[] An assessment of the likely outcome(s) for a patient or other subject as well as the likelihood of each outcome.
Method specifies the algorithm, process or mechanism used to evaluate the risk.
Code specifies the type of the risk assessment performed.
For example, in the system  \url{http://snomed.info/sct}, `709510001' represents Assessment of risk for disease (procedure).
For assessments or prognosis specific to a particular condition, the condition attribute indicates the condition being assessed. It should be a reference to an existing condition but is a simplified as a string here.
Reason specifies the reason the risk assessment was performed.
Mitigation contains a description of the steps that might be taken to reduce the identified risk(s).
\end{itemize}

\item \vpedatatype{data:Datatypesorghl7fhirr5modelRiskAssessmentRiskAssessmentPredictionComponent}{RiskAssessmentPredictionComponent}: 
\begin{itemize}[noitemsep,nolistsep]
\item[] Attributes: String outcome, \ref{data:Datatypesorghl7fhirr5modelRange} probability, String qualitativeRisk, \ref{data:DatatypesBigDecimal} relativeRisk, \ref{data:DatatypesDate} when, String rationale, long serialVersionUID
\item[] Describes the expected outcome for the subject.
Outcome speciifes one of the potential outcomes for the patient (e.g. remission, death,  a particular condition).
QualitativeRisk indicates how likely the outcome is (in the specified timeframe), expressed as a qualitative value (e.g. low, medium, or high).
Rationale contains additional information explaining the basis for the prediction.
Both outcome and qualitativeRisk are CodeableConcepts (which can be formal reference to terminology or an ontology), but simplified as a string.
\end{itemize}

\item \vpedatatype{data:DatatypesRiskEstimationId}{RiskEstimationId}: 
\begin{itemize}[noitemsep,nolistsep]

\item[] A piece of data uniquely identifying a risk estimation performed by the PMS. This architecture does not specify the exact format of this identifier,
but possibilities are a long integer, a string, a URL etc.
\end{itemize}

\item \vpedatatype{data:DatatypesExceptionsSecurityException}{\textit{SecurityException}}: 
\begin{itemize}[noitemsep,nolistsep]

\item[] Parent class for security exceptions.
\end{itemize}

\item \vpedatatype{data:DatatypesSensorDataElement}{SensorDataElement}: 
\begin{itemize}[noitemsep,nolistsep]
\item[] Attributes: String sensorId, String sensorType, String measurementType, String measurementUnit, \ref{data:DatatypesBigDecimal} measurementValue, String measurementValueString
\item[] Each sub-package lists the id of the sensor, the type of the sensor, the
type of the measurements and the measurements itself. These
measurements range from a single value (e.g., in case of the blood
pressure) to a complex data structure (e.g., in case of an ECG).
Non-numerical types of measurements can be captured in the measurementValueString attribute.
\end{itemize}

\item \vpedatatype{data:DatatypesSensorDataPackage}{SensorDataPackage}: 
\begin{itemize}[noitemsep,nolistsep]
\item[] Attributes: List\textless{}\ref{data:DatatypesSensorDataElement}\textgreater{} sensorDataElements
\item[] A package of sensor data, i.e., a list of sensorDataElements.

\end{itemize}

\item \vpedatatype{data:string}{string}: 
\begin{itemize}[noitemsep,nolistsep]

\item[] {\colorbox{red!30}{\underline{Undefined}}}
\end{itemize}

\item \vpedatatype{data:DatatypesThirdPartyDataElement}{ThirdPartyDataElement}: 
\begin{itemize}[noitemsep,nolistsep]
\item[] Attributes: String serviceType, String measurementType, String measurementUnit, \ref{data:DatatypesBigDecimal} measurementValue, String measurementString
\item[] Third party data element is a struct which can store any type obtained from the third party fitness tracker
\end{itemize}

\item \vpedatatype{data:DatatypesThirdPartyDataPackage}{ThirdPartyDataPackage}: 
\begin{itemize}[noitemsep,nolistsep]
\item[] Attributes: List\textless{}\ref{data:DatatypesThirdPartyDataElement}\textgreater{} thirdPartyDataElements
\item[] class for representing the data polled from the third party fitness trackers
\end{itemize}

\item \vpedatatype{data:DatatypesTimestamp}{Timestamp}: 
\begin{itemize}[noitemsep,nolistsep]

\item[] The representation of a time (i.e., date and time of the day) in the system.
\end{itemize}

\item \vpedatatype{data:DatatypesTokenPackage}{TokenPackage}: 
\begin{itemize}[noitemsep,nolistsep]
\item[] Attributes: \ref{data:DatatypesPatientId} patientId, string token, string provider
\item[] a package to store the important info of a token used to connect to a third party app api.  
\end{itemize}

\end{itemize}
% END TYPES

% CHECKS
\section{Unresolved issues}\label{sec:checks}
\emph{SA Plugin v10.0.0 (VP OpenAPI v17.2)}
\begin{itemize}[nolistsep,noitemsep]
\item[] No failed checks
\end{itemize}
% END CHECKS

%\end{document}
